% paper55-trust-economics.tex — Economic Models for Trust: Verification as a Market
%   Paper 55 of the Judgment Geometry series.
% Compile: pdflatex paper55-trust-economics
\documentclass[11pt]{article}
\usepackage{lmodern}
% jugeo-common.tex — shared preamble for all Judgment Geometry papers
% Include via: % jugeo-common.tex — shared preamble for all Judgment Geometry papers
% Include via: % jugeo-common.tex — shared preamble for all Judgment Geometry papers
% Include via: \input{jugeo-common}

% ── Packages ────────────────────────────────────────────────────────────
\usepackage{amsmath,amssymb,amsthm,mathtools}
\usepackage{tikz-cd}
\usepackage{listings}
\usepackage{xcolor}
\usepackage{xspace}
\usepackage{booktabs}
\usepackage{multirow}
\usepackage{enumitem}
\usepackage[expansion=false]{microtype}
\usepackage{hyperref}
\usepackage{cleveref}
% \usepackage{thmtools}  % disabled: not available in this TeX installation
\usepackage{float}
\usepackage{caption}
\usepackage{subcaption}
\usepackage{tabularx}
\usepackage{stmaryrd}       % \llbracket, \rrbracket
\usepackage{bbm}            % \mathbbm{1}

% ── Page geometry ───────────────────────────────────────────────────────
\usepackage[margin=1in]{geometry}

% ── Theorem environments ────────────────────────────────────────────────
\theoremstyle{plain}
\newtheorem{theorem}{Theorem}[section]
\newtheorem{lemma}[theorem]{Lemma}
\newtheorem{proposition}[theorem]{Proposition}
\newtheorem{corollary}[theorem]{Corollary}

\theoremstyle{definition}
\newtheorem{definition}[theorem]{Definition}
\newtheorem{example}[theorem]{Example}
\newtheorem{construction}[theorem]{Construction}

\theoremstyle{remark}
\newtheorem{remark}[theorem]{Remark}
\newtheorem{notation}[theorem]{Notation}

% ── Python listings style ──────────────────────────────────────────────
\definecolor{jugeoBlue}{HTML}{2563EB}
\definecolor{jugeoGreen}{HTML}{059669}
\definecolor{jugeoRed}{HTML}{DC2626}
\definecolor{jugeoPurple}{HTML}{7C3AED}
\definecolor{jugeoGray}{HTML}{6B7280}
\definecolor{jugeoBg}{HTML}{F8FAFC}

\lstdefinestyle{jugeo-python}{
  language=Python,
  basicstyle=\ttfamily\small,
  keywordstyle=\color{jugeoBlue}\bfseries,
  stringstyle=\color{jugeoGreen},
  commentstyle=\color{jugeoGray}\itshape,
  numberstyle=\tiny\color{jugeoGray},
  backgroundcolor=\color{jugeoBg},
  numbers=left,
  numbersep=6pt,
  frame=single,
  framerule=0.4pt,
  rulecolor=\color{jugeoGray!40},
  breaklines=true,
  breakatwhitespace=true,
  showstringspaces=false,
  tabsize=4,
  xleftmargin=1.5em,
  framexleftmargin=1.5em,
  morekeywords={dataclass,frozen,field,Enum,IntEnum,auto,Any,
                Mapping,Sequence,Optional,match,case},
  literate={→}{{$\to$}}1 {←}{{$\leftarrow$}}1
           {≤}{{$\leq$}}1 {≥}{{$\geq$}}1 {≠}{{$\neq$}}1
           {∈}{{$\in$}}1 {∀}{{$\forall$}}1 {∃}{{$\exists$}}1
           {⊕}{{$\oplus$}}1 {⊖}{{$\ominus$}}1,
  escapeinside={(*@}{@*)},
}
\lstset{style=jugeo-python}

\lstdefinestyle{jugeo-output}{
  basicstyle=\ttfamily\small,
  backgroundcolor=\color{jugeoBg},
  frame=single,
  framerule=0.4pt,
  rulecolor=\color{jugeoGray!40},
  breaklines=true,
  showstringspaces=false,
  xleftmargin=1.5em,
  framexleftmargin=0.5em,
  numbers=none,
}

% ── JuGeo-specific macros ──────────────────────────────────────────────

% Judgment 8-tuple
\newcommand{\J}{\mathcal{J}}
\newcommand{\Jud}[8]{(#1,\,#2,\,#3,\,#4,\,#5,\,#6,\,#7,\,#8)}
\newcommand{\JudShort}[1]{\J_{#1}}

% Components
\newcommand{\coord}{\mathsf{c}}            % coordinate
\newcommand{\prop}{\varphi}                 % proposition
\newcommand{\carrier}{\mathcal{A}}          % carrier type
\newcommand{\evid}{\mathcal{E}}             % evidence bundle
\newcommand{\oblig}{\mathcal{O}}            % obligations
\newcommand{\obstr}{\mathcal{B}}            % obstructions
\newcommand{\trust}{\mathcal{T}}            % trust annotation
\newcommand{\prov}{\Pi}                     % provenance

% Trust algebra
\newcommand{\Talg}{\mathfrak{T}}            % trust ordered algebra
\newcommand{\Eadm}{\mathcal{E}_{\mathrm{adm}}}
\newcommand{\tleq}{\preceq}                 % trust ordering
\newcommand{\tjoin}{\oplus}                 % conservative join
\newcommand{\tatten}{\ominus}               % attenuation
\newcommand{\tpromote}{\uparrow_\pi}        % promotion
\newcommand{\tdemote}{\downarrow_\chi}       % demotion

% Trust tiers
\newcommand{\Tcontra}{\ensuremath{\mathsf{CONTRADICTED}}}
\newcommand{\Tunver}{\ensuremath{\mathsf{UNVERIFIED}}}
\newcommand{\Tcopilot}{\ensuremath{\mathsf{COPILOT}}}
\newcommand{\Toracle}{\ensuremath{\mathsf{ORACLE}}}
\newcommand{\Truntime}{\ensuremath{\mathsf{RUNTIME}}}
\newcommand{\Tsolver}{\ensuremath{\mathsf{SOLVER}}}
\newcommand{\Tproof}{\ensuremath{\mathsf{PROOF}}}

% Site and geometry
\newcommand{\Site}{\mathcal{S}}             % semantic site
\newcommand{\Cov}{\mathrm{Cov}}            % covering family
\newcommand{\Ob}{\mathrm{Ob}}              % objects
\newcommand{\Mor}{\mathrm{Mor}}            % morphisms
\newcommand{\res}{\mathrm{res}}            % restriction
\newcommand{\Cech}{\check{C}}              % Čech complex
\newcommand{\Hcoh}[1]{H^{#1}}             % cohomology group

% Descent
\newcommand{\descent}{\mathrm{desc}}
\newcommand{\glue}{\mathrm{glue}}
\newcommand{\overlap}[2]{#1 \cap #2}

% Semantic moves
\newcommand{\VERIFY}{\textsc{Verify}}
\newcommand{\CONSTRUCT}{\textsc{Construct}}
\newcommand{\NORMALIZE}{\textsc{Normalize}}
\newcommand{\GLUE}{\textsc{Glue}}
\newcommand{\RESTRICT}{\textsc{Restrict}}
\newcommand{\TRANSPORT}{\textsc{Transport}}
\newcommand{\REPAIR}{\textsc{Repair}}
\newcommand{\PROMOTE}{\textsc{Promote}}
\newcommand{\CHALLENGE}{\textsc{Challenge}}
\newcommand{\ESCALATE}{\textsc{Escalate}}

% Fragments
\newcommand{\QFLIA}{\texttt{QF\_LIA}}
\newcommand{\QFLRA}{\texttt{QF\_LRA}}
\newcommand{\QFBV}{\texttt{QF\_BV}}
\newcommand{\QFUF}{\texttt{QF\_UF}}

% Misc
\newcommand{\Python}{\textsc{Python}}
\newcommand{\JuGeo}{\textsc{JuGeo}}
\newcommand{\LEAN}{\textsc{Lean}\,4}
\newcommand{\Fstar}{F$^{\star}$}
\newcommand{\Coq}{\textsc{Coq}}
\newcommand{\Dafny}{\textsc{Dafny}}

% Common shortcuts
\newcommand{\ie}{i.e.\@\xspace}
\newcommand{\eg}{e.g.\@\xspace}
\newcommand{\cf}{cf.\@\xspace}
\newcommand{\wrt}{w.r.t.\@\xspace}

% Evaluation shortcuts
\newcommand{\numcases}{300}
\newcommand{\numspec}{100}
\newcommand{\numeq}{100}
\newcommand{\numbug}{100}
\newcommand{\medtime}{6.5\,ms}
\newcommand{\totaltime}{1.949\,s}


% ── Packages ────────────────────────────────────────────────────────────
\usepackage{amsmath,amssymb,amsthm,mathtools}
\usepackage{tikz-cd}
\usepackage{listings}
\usepackage{xcolor}
\usepackage{xspace}
\usepackage{booktabs}
\usepackage{multirow}
\usepackage{enumitem}
\usepackage[expansion=false]{microtype}
\usepackage{hyperref}
\usepackage{cleveref}
% \usepackage{thmtools}  % disabled: not available in this TeX installation
\usepackage{float}
\usepackage{caption}
\usepackage{subcaption}
\usepackage{tabularx}
\usepackage{stmaryrd}       % \llbracket, \rrbracket
\usepackage{bbm}            % \mathbbm{1}

% ── Page geometry ───────────────────────────────────────────────────────
\usepackage[margin=1in]{geometry}

% ── Theorem environments ────────────────────────────────────────────────
\theoremstyle{plain}
\newtheorem{theorem}{Theorem}[section]
\newtheorem{lemma}[theorem]{Lemma}
\newtheorem{proposition}[theorem]{Proposition}
\newtheorem{corollary}[theorem]{Corollary}

\theoremstyle{definition}
\newtheorem{definition}[theorem]{Definition}
\newtheorem{example}[theorem]{Example}
\newtheorem{construction}[theorem]{Construction}

\theoremstyle{remark}
\newtheorem{remark}[theorem]{Remark}
\newtheorem{notation}[theorem]{Notation}

% ── Python listings style ──────────────────────────────────────────────
\definecolor{jugeoBlue}{HTML}{2563EB}
\definecolor{jugeoGreen}{HTML}{059669}
\definecolor{jugeoRed}{HTML}{DC2626}
\definecolor{jugeoPurple}{HTML}{7C3AED}
\definecolor{jugeoGray}{HTML}{6B7280}
\definecolor{jugeoBg}{HTML}{F8FAFC}

\lstdefinestyle{jugeo-python}{
  language=Python,
  basicstyle=\ttfamily\small,
  keywordstyle=\color{jugeoBlue}\bfseries,
  stringstyle=\color{jugeoGreen},
  commentstyle=\color{jugeoGray}\itshape,
  numberstyle=\tiny\color{jugeoGray},
  backgroundcolor=\color{jugeoBg},
  numbers=left,
  numbersep=6pt,
  frame=single,
  framerule=0.4pt,
  rulecolor=\color{jugeoGray!40},
  breaklines=true,
  breakatwhitespace=true,
  showstringspaces=false,
  tabsize=4,
  xleftmargin=1.5em,
  framexleftmargin=1.5em,
  morekeywords={dataclass,frozen,field,Enum,IntEnum,auto,Any,
                Mapping,Sequence,Optional,match,case},
  literate={→}{{$\to$}}1 {←}{{$\leftarrow$}}1
           {≤}{{$\leq$}}1 {≥}{{$\geq$}}1 {≠}{{$\neq$}}1
           {∈}{{$\in$}}1 {∀}{{$\forall$}}1 {∃}{{$\exists$}}1
           {⊕}{{$\oplus$}}1 {⊖}{{$\ominus$}}1,
  escapeinside={(*@}{@*)},
}
\lstset{style=jugeo-python}

\lstdefinestyle{jugeo-output}{
  basicstyle=\ttfamily\small,
  backgroundcolor=\color{jugeoBg},
  frame=single,
  framerule=0.4pt,
  rulecolor=\color{jugeoGray!40},
  breaklines=true,
  showstringspaces=false,
  xleftmargin=1.5em,
  framexleftmargin=0.5em,
  numbers=none,
}

% ── JuGeo-specific macros ──────────────────────────────────────────────

% Judgment 8-tuple
\newcommand{\J}{\mathcal{J}}
\newcommand{\Jud}[8]{(#1,\,#2,\,#3,\,#4,\,#5,\,#6,\,#7,\,#8)}
\newcommand{\JudShort}[1]{\J_{#1}}

% Components
\newcommand{\coord}{\mathsf{c}}            % coordinate
\newcommand{\prop}{\varphi}                 % proposition
\newcommand{\carrier}{\mathcal{A}}          % carrier type
\newcommand{\evid}{\mathcal{E}}             % evidence bundle
\newcommand{\oblig}{\mathcal{O}}            % obligations
\newcommand{\obstr}{\mathcal{B}}            % obstructions
\newcommand{\trust}{\mathcal{T}}            % trust annotation
\newcommand{\prov}{\Pi}                     % provenance

% Trust algebra
\newcommand{\Talg}{\mathfrak{T}}            % trust ordered algebra
\newcommand{\Eadm}{\mathcal{E}_{\mathrm{adm}}}
\newcommand{\tleq}{\preceq}                 % trust ordering
\newcommand{\tjoin}{\oplus}                 % conservative join
\newcommand{\tatten}{\ominus}               % attenuation
\newcommand{\tpromote}{\uparrow_\pi}        % promotion
\newcommand{\tdemote}{\downarrow_\chi}       % demotion

% Trust tiers
\newcommand{\Tcontra}{\ensuremath{\mathsf{CONTRADICTED}}}
\newcommand{\Tunver}{\ensuremath{\mathsf{UNVERIFIED}}}
\newcommand{\Tcopilot}{\ensuremath{\mathsf{COPILOT}}}
\newcommand{\Toracle}{\ensuremath{\mathsf{ORACLE}}}
\newcommand{\Truntime}{\ensuremath{\mathsf{RUNTIME}}}
\newcommand{\Tsolver}{\ensuremath{\mathsf{SOLVER}}}
\newcommand{\Tproof}{\ensuremath{\mathsf{PROOF}}}

% Site and geometry
\newcommand{\Site}{\mathcal{S}}             % semantic site
\newcommand{\Cov}{\mathrm{Cov}}            % covering family
\newcommand{\Ob}{\mathrm{Ob}}              % objects
\newcommand{\Mor}{\mathrm{Mor}}            % morphisms
\newcommand{\res}{\mathrm{res}}            % restriction
\newcommand{\Cech}{\check{C}}              % Čech complex
\newcommand{\Hcoh}[1]{H^{#1}}             % cohomology group

% Descent
\newcommand{\descent}{\mathrm{desc}}
\newcommand{\glue}{\mathrm{glue}}
\newcommand{\overlap}[2]{#1 \cap #2}

% Semantic moves
\newcommand{\VERIFY}{\textsc{Verify}}
\newcommand{\CONSTRUCT}{\textsc{Construct}}
\newcommand{\NORMALIZE}{\textsc{Normalize}}
\newcommand{\GLUE}{\textsc{Glue}}
\newcommand{\RESTRICT}{\textsc{Restrict}}
\newcommand{\TRANSPORT}{\textsc{Transport}}
\newcommand{\REPAIR}{\textsc{Repair}}
\newcommand{\PROMOTE}{\textsc{Promote}}
\newcommand{\CHALLENGE}{\textsc{Challenge}}
\newcommand{\ESCALATE}{\textsc{Escalate}}

% Fragments
\newcommand{\QFLIA}{\texttt{QF\_LIA}}
\newcommand{\QFLRA}{\texttt{QF\_LRA}}
\newcommand{\QFBV}{\texttt{QF\_BV}}
\newcommand{\QFUF}{\texttt{QF\_UF}}

% Misc
\newcommand{\Python}{\textsc{Python}}
\newcommand{\JuGeo}{\textsc{JuGeo}}
\newcommand{\LEAN}{\textsc{Lean}\,4}
\newcommand{\Fstar}{F$^{\star}$}
\newcommand{\Coq}{\textsc{Coq}}
\newcommand{\Dafny}{\textsc{Dafny}}

% Common shortcuts
\newcommand{\ie}{i.e.\@\xspace}
\newcommand{\eg}{e.g.\@\xspace}
\newcommand{\cf}{cf.\@\xspace}
\newcommand{\wrt}{w.r.t.\@\xspace}

% Evaluation shortcuts
\newcommand{\numcases}{300}
\newcommand{\numspec}{100}
\newcommand{\numeq}{100}
\newcommand{\numbug}{100}
\newcommand{\medtime}{6.5\,ms}
\newcommand{\totaltime}{1.949\,s}


% ── Packages ────────────────────────────────────────────────────────────
\usepackage{amsmath,amssymb,amsthm,mathtools}
\usepackage{tikz-cd}
\usepackage{listings}
\usepackage{xcolor}
\usepackage{xspace}
\usepackage{booktabs}
\usepackage{multirow}
\usepackage{enumitem}
\usepackage[expansion=false]{microtype}
\usepackage{hyperref}
\usepackage{cleveref}
% \usepackage{thmtools}  % disabled: not available in this TeX installation
\usepackage{float}
\usepackage{caption}
\usepackage{subcaption}
\usepackage{tabularx}
\usepackage{stmaryrd}       % \llbracket, \rrbracket
\usepackage{bbm}            % \mathbbm{1}

% ── Page geometry ───────────────────────────────────────────────────────
\usepackage[margin=1in]{geometry}

% ── Theorem environments ────────────────────────────────────────────────
\theoremstyle{plain}
\newtheorem{theorem}{Theorem}[section]
\newtheorem{lemma}[theorem]{Lemma}
\newtheorem{proposition}[theorem]{Proposition}
\newtheorem{corollary}[theorem]{Corollary}

\theoremstyle{definition}
\newtheorem{definition}[theorem]{Definition}
\newtheorem{example}[theorem]{Example}
\newtheorem{construction}[theorem]{Construction}

\theoremstyle{remark}
\newtheorem{remark}[theorem]{Remark}
\newtheorem{notation}[theorem]{Notation}

% ── Python listings style ──────────────────────────────────────────────
\definecolor{jugeoBlue}{HTML}{2563EB}
\definecolor{jugeoGreen}{HTML}{059669}
\definecolor{jugeoRed}{HTML}{DC2626}
\definecolor{jugeoPurple}{HTML}{7C3AED}
\definecolor{jugeoGray}{HTML}{6B7280}
\definecolor{jugeoBg}{HTML}{F8FAFC}

\lstdefinestyle{jugeo-python}{
  language=Python,
  basicstyle=\ttfamily\small,
  keywordstyle=\color{jugeoBlue}\bfseries,
  stringstyle=\color{jugeoGreen},
  commentstyle=\color{jugeoGray}\itshape,
  numberstyle=\tiny\color{jugeoGray},
  backgroundcolor=\color{jugeoBg},
  numbers=left,
  numbersep=6pt,
  frame=single,
  framerule=0.4pt,
  rulecolor=\color{jugeoGray!40},
  breaklines=true,
  breakatwhitespace=true,
  showstringspaces=false,
  tabsize=4,
  xleftmargin=1.5em,
  framexleftmargin=1.5em,
  morekeywords={dataclass,frozen,field,Enum,IntEnum,auto,Any,
                Mapping,Sequence,Optional,match,case},
  literate={→}{{$\to$}}1 {←}{{$\leftarrow$}}1
           {≤}{{$\leq$}}1 {≥}{{$\geq$}}1 {≠}{{$\neq$}}1
           {∈}{{$\in$}}1 {∀}{{$\forall$}}1 {∃}{{$\exists$}}1
           {⊕}{{$\oplus$}}1 {⊖}{{$\ominus$}}1,
  escapeinside={(*@}{@*)},
}
\lstset{style=jugeo-python}

\lstdefinestyle{jugeo-output}{
  basicstyle=\ttfamily\small,
  backgroundcolor=\color{jugeoBg},
  frame=single,
  framerule=0.4pt,
  rulecolor=\color{jugeoGray!40},
  breaklines=true,
  showstringspaces=false,
  xleftmargin=1.5em,
  framexleftmargin=0.5em,
  numbers=none,
}

% ── JuGeo-specific macros ──────────────────────────────────────────────

% Judgment 8-tuple
\newcommand{\J}{\mathcal{J}}
\newcommand{\Jud}[8]{(#1,\,#2,\,#3,\,#4,\,#5,\,#6,\,#7,\,#8)}
\newcommand{\JudShort}[1]{\J_{#1}}

% Components
\newcommand{\coord}{\mathsf{c}}            % coordinate
\newcommand{\prop}{\varphi}                 % proposition
\newcommand{\carrier}{\mathcal{A}}          % carrier type
\newcommand{\evid}{\mathcal{E}}             % evidence bundle
\newcommand{\oblig}{\mathcal{O}}            % obligations
\newcommand{\obstr}{\mathcal{B}}            % obstructions
\newcommand{\trust}{\mathcal{T}}            % trust annotation
\newcommand{\prov}{\Pi}                     % provenance

% Trust algebra
\newcommand{\Talg}{\mathfrak{T}}            % trust ordered algebra
\newcommand{\Eadm}{\mathcal{E}_{\mathrm{adm}}}
\newcommand{\tleq}{\preceq}                 % trust ordering
\newcommand{\tjoin}{\oplus}                 % conservative join
\newcommand{\tatten}{\ominus}               % attenuation
\newcommand{\tpromote}{\uparrow_\pi}        % promotion
\newcommand{\tdemote}{\downarrow_\chi}       % demotion

% Trust tiers
\newcommand{\Tcontra}{\ensuremath{\mathsf{CONTRADICTED}}}
\newcommand{\Tunver}{\ensuremath{\mathsf{UNVERIFIED}}}
\newcommand{\Tcopilot}{\ensuremath{\mathsf{COPILOT}}}
\newcommand{\Toracle}{\ensuremath{\mathsf{ORACLE}}}
\newcommand{\Truntime}{\ensuremath{\mathsf{RUNTIME}}}
\newcommand{\Tsolver}{\ensuremath{\mathsf{SOLVER}}}
\newcommand{\Tproof}{\ensuremath{\mathsf{PROOF}}}

% Site and geometry
\newcommand{\Site}{\mathcal{S}}             % semantic site
\newcommand{\Cov}{\mathrm{Cov}}            % covering family
\newcommand{\Ob}{\mathrm{Ob}}              % objects
\newcommand{\Mor}{\mathrm{Mor}}            % morphisms
\newcommand{\res}{\mathrm{res}}            % restriction
\newcommand{\Cech}{\check{C}}              % Čech complex
\newcommand{\Hcoh}[1]{H^{#1}}             % cohomology group

% Descent
\newcommand{\descent}{\mathrm{desc}}
\newcommand{\glue}{\mathrm{glue}}
\newcommand{\overlap}[2]{#1 \cap #2}

% Semantic moves
\newcommand{\VERIFY}{\textsc{Verify}}
\newcommand{\CONSTRUCT}{\textsc{Construct}}
\newcommand{\NORMALIZE}{\textsc{Normalize}}
\newcommand{\GLUE}{\textsc{Glue}}
\newcommand{\RESTRICT}{\textsc{Restrict}}
\newcommand{\TRANSPORT}{\textsc{Transport}}
\newcommand{\REPAIR}{\textsc{Repair}}
\newcommand{\PROMOTE}{\textsc{Promote}}
\newcommand{\CHALLENGE}{\textsc{Challenge}}
\newcommand{\ESCALATE}{\textsc{Escalate}}

% Fragments
\newcommand{\QFLIA}{\texttt{QF\_LIA}}
\newcommand{\QFLRA}{\texttt{QF\_LRA}}
\newcommand{\QFBV}{\texttt{QF\_BV}}
\newcommand{\QFUF}{\texttt{QF\_UF}}

% Misc
\newcommand{\Python}{\textsc{Python}}
\newcommand{\JuGeo}{\textsc{JuGeo}}
\newcommand{\LEAN}{\textsc{Lean}\,4}
\newcommand{\Fstar}{F$^{\star}$}
\newcommand{\Coq}{\textsc{Coq}}
\newcommand{\Dafny}{\textsc{Dafny}}

% Common shortcuts
\newcommand{\ie}{i.e.\@\xspace}
\newcommand{\eg}{e.g.\@\xspace}
\newcommand{\cf}{cf.\@\xspace}
\newcommand{\wrt}{w.r.t.\@\xspace}

% Evaluation shortcuts
\newcommand{\numcases}{300}
\newcommand{\numspec}{100}
\newcommand{\numeq}{100}
\newcommand{\numbug}{100}
\newcommand{\medtime}{6.5\,ms}
\newcommand{\totaltime}{1.949\,s}

% experiment-data.tex — AUTO-GENERATED from real jugeo CLI runs
% DO NOT EDIT — regenerate with: python3 experiments/run_universal_metrics.py
% Generated from 30 programs across 6 domains

\newcommand{\expTotalPrograms}{30}
\newcommand{\expVerified}{30}
\newcommand{\expAccuracy}{100.0\%}
\newcommand{\expCoordsMin}{2}
\newcommand{\expCoordsMax}{10}
\newcommand{\expCoordsMean}{5.2}
\newcommand{\expCoordsMedian}{5.5}
\newcommand{\expPropsMin}{4}
\newcommand{\expPropsMax}{20}
\newcommand{\expPropsMean}{10.4}
\newcommand{\expPropsSum}{311}
\newcommand{\expPropsOkSum}{310}
\newcommand{\expTimeMin}{3.506\,s}
\newcommand{\expTimeMax}{6.714\,s}
\newcommand{\expTimeMean}{4.64\,s}
\newcommand{\expTimeMedian}{4.508\,s}
\newcommand{\expTimeTotal}{139.204\,s}
\newcommand{\expObstructionTotal}{0}
\newcommand{\expHOneZero}{30}
\newcommand{\expDomainCount}{6}
\newcommand{\expTrustCOPILOTSUGGESTED}{30}
\newcommand{\expDomainDsCount}{5}
\newcommand{\expDomainDsVerified}{5}
\newcommand{\expDomainDsAvgCoords}{7.6}
\newcommand{\expDomainDsAvgProps}{15.2}
\newcommand{\expDomainMathCount}{5}
\newcommand{\expDomainMathVerified}{5}
\newcommand{\expDomainMathAvgCoords}{4.8}
\newcommand{\expDomainMathAvgProps}{9.4}
\newcommand{\expDomainSmCount}{5}
\newcommand{\expDomainSmVerified}{5}
\newcommand{\expDomainSmAvgCoords}{7.6}
\newcommand{\expDomainSmAvgProps}{15.2}
\newcommand{\expDomainSortCount}{5}
\newcommand{\expDomainSortVerified}{5}
\newcommand{\expDomainSortAvgCoords}{2.4}
\newcommand{\expDomainSortAvgProps}{4.8}
\newcommand{\expDomainStrCount}{5}
\newcommand{\expDomainStrVerified}{5}
\newcommand{\expDomainStrAvgCoords}{3.2}
\newcommand{\expDomainStrAvgProps}{6.4}
\newcommand{\expDomainWebCount}{5}
\newcommand{\expDomainWebVerified}{5}
\newcommand{\expDomainWebAvgCoords}{5.6}
\newcommand{\expDomainWebAvgProps}{11.2}

% subsystem-data.tex — AUTO-GENERATED from real jugeo subsystem experiments
% DO NOT EDIT — regenerate with: python3 experiments/run_subsystem_experiments.py

\newcommand{\subCliTotal}{10}
\newcommand{\subCliVerified}{9}
\newcommand{\subCliAccuracy}{90.0\%}
\newcommand{\subCliMeanTime}{4.132\,s}
\newcommand{\subCliMinTime}{3.472\,s}
\newcommand{\subCliMaxTime}{5.372\,s}
\newcommand{\subCliTotalCoords}{54}
\newcommand{\subCliTotalProps}{107}
\newcommand{\subCliTotalPropsOk}{106}
\newcommand{\subSiteCalls}{90}
\newcommand{\subSiteSuccessful}{69}
\newcommand{\subSiteSuccessRate}{76.7\%}
\newcommand{\subSiteMeanTime}{0.0001\,s}
\newcommand{\subSiteTotalTime}{0.005\,s}
\newcommand{\subSpecificationsatisfactionSmallTime}{0.0003\,s}
\newcommand{\subSpecificationsatisfactionMediumTime}{0.0001\,s}
\newcommand{\subSpecificationsatisfactionLargeTime}{0.0001\,s}
\newcommand{\subInhabitantfleetSmallTime}{0.0\,s}
\newcommand{\subInhabitantfleetMediumTime}{0.0\,s}
\newcommand{\subInhabitantfleetLargeTime}{0.0\,s}
\newcommand{\subTheoremecologySmallTime}{0.0\,s}
\newcommand{\subTheoremecologyMediumTime}{0.0\,s}
\newcommand{\subTheoremecologyLargeTime}{0.0\,s}
\newcommand{\subStatespaceexplorationSmallTime}{0.0\,s}
\newcommand{\subStatespaceexplorationMediumTime}{0.0\,s}
\newcommand{\subStatespaceexplorationLargeTime}{0.0\,s}
\newcommand{\subRepairsemanticsSmallTime}{0.0\,s}
\newcommand{\subRepairsemanticsMediumTime}{0.0\,s}
\newcommand{\subRepairsemanticsLargeTime}{0.0\,s}
\newcommand{\subReplaygluingSmallTime}{0.0\,s}
\newcommand{\subReplaygluingMediumTime}{0.0\,s}
\newcommand{\subReplaygluingLargeTime}{0.0\,s}
\newcommand{\subSemanticfuturesSmallTime}{0.0\,s}
\newcommand{\subSemanticfuturesMediumTime}{0.0\,s}
\newcommand{\subSemanticfuturesLargeTime}{0.0\,s}
\newcommand{\subHypercovertreatySmallTime}{0.0\,s}
\newcommand{\subHypercovertreatyMediumTime}{0.0\,s}
\newcommand{\subHypercovertreatyLargeTime}{0.0\,s}
\newcommand{\subEncodeforsolverSmallTime}{0.0026\,s}
\newcommand{\subEncodeforsolverMediumTime}{0.0002\,s}
\newcommand{\subEncodeforsolverLargeTime}{0.0001\,s}
\newcommand{\subInterfaceroutingSmallTime}{0.0\,s}
\newcommand{\subInterfaceroutingMediumTime}{0.0\,s}
\newcommand{\subInterfaceroutingLargeTime}{0.0\,s}
\newcommand{\subOrchestrateverificationSmallTime}{0.0002\,s}
\newcommand{\subOrchestrateverificationMediumTime}{0.0\,s}
\newcommand{\subOrchestrateverificationLargeTime}{0.0\,s}
\newcommand{\subRunfulldescentSmallTime}{0.0\,s}
\newcommand{\subRunfulldescentMediumTime}{0.0\,s}
\newcommand{\subRunfulldescentLargeTime}{0.0\,s}
\newcommand{\subKernellifecycleSmallTime}{0.0\,s}
\newcommand{\subKernellifecycleMediumTime}{0.0\,s}
\newcommand{\subKernellifecycleLargeTime}{0.0\,s}
\newcommand{\subDiscoverypipelineSmallTime}{0.0\,s}
\newcommand{\subDiscoverypipelineMediumTime}{0.0\,s}
\newcommand{\subDiscoverypipelineLargeTime}{0.0\,s}
\newcommand{\subAnalogytransportSmallTime}{0.0\,s}
\newcommand{\subAnalogytransportMediumTime}{0.0\,s}
\newcommand{\subAnalogytransportLargeTime}{0.0\,s}
\newcommand{\subFormalcoresiteSmallTime}{0.0001\,s}
\newcommand{\subFormalcoresiteMediumTime}{0.0\,s}
\newcommand{\subFormalcoresiteLargeTime}{0.0\,s}
\newcommand{\subProblematlasSmallTime}{0.0\,s}
\newcommand{\subProblematlasMediumTime}{0.0\,s}
\newcommand{\subProblematlasLargeTime}{0.0\,s}
\newcommand{\subPublicalignmentSmallTime}{0.0\,s}
\newcommand{\subPublicalignmentMediumTime}{0.0\,s}
\newcommand{\subPublicalignmentLargeTime}{0.0\,s}
\newcommand{\subRegimebootstrappingSmallTime}{0.0\,s}
\newcommand{\subRegimebootstrappingMediumTime}{0.0\,s}
\newcommand{\subRegimebootstrappingLargeTime}{0.0\,s}
\newcommand{\subRelationalrefinementSmallTime}{0.0\,s}
\newcommand{\subRelationalrefinementMediumTime}{0.0\,s}
\newcommand{\subRelationalrefinementLargeTime}{0.0\,s}
\newcommand{\subSemanticclosureSmallTime}{0.0\,s}
\newcommand{\subSemanticclosureMediumTime}{0.0\,s}
\newcommand{\subSemanticclosureLargeTime}{0.0\,s}
\newcommand{\subTheoremeconomicsSmallTime}{0.0001\,s}
\newcommand{\subTheoremeconomicsMediumTime}{0.0\,s}
\newcommand{\subTheoremeconomicsLargeTime}{0.0\,s}
\newcommand{\subBenchmarksuiteSmallTime}{0.0\,s}
\newcommand{\subBenchmarksuiteMediumTime}{0.0\,s}
\newcommand{\subBenchmarksuiteLargeTime}{0.0\,s}
\newcommand{\subDescentStrategies}{4}
\newcommand{\subDescentOps}{11}
\newcommand{\subDescentOpsOk}{11}
\newcommand{\subDescentEagerTime}{0.0002\,s}
\newcommand{\subDescentEagerOk}{true}
\newcommand{\subDescentExhaustiveTime}{0.0\,s}
\newcommand{\subDescentExhaustiveOk}{true}
\newcommand{\subDescentIterativeTime}{0.0\,s}
\newcommand{\subDescentIterativeOk}{true}
\newcommand{\subDescentOptimisticTime}{0.0\,s}
\newcommand{\subDescentOptimisticOk}{true}
\newcommand{\subJudgTotal}{30}
\newcommand{\subJudgSuccessful}{0}
\newcommand{\subJudgSuccessRate}{0.0\%}
\newcommand{\subJudgMeanTimeUs}{7.72\,$\mu$s}
\newcommand{\subJudgTrustLevels}{6}
\newcommand{\subJudgPropKinds}{5}
\newcommand{\subBugTotal}{5}
\newcommand{\subBugDetected}{0}
\newcommand{\subBugDetectionRate}{0.0\%}
\newcommand{\subBugFalsePositiveRate}{0.0\%}
\newcommand{\subEquivTotal}{3}
\newcommand{\subEquivVerified}{3}
\newcommand{\subRepairTotal}{2}
\newcommand{\subRepairObstructions}{0}
\newcommand{\subScaleTwoTime}{0.0001\,s}
\newcommand{\subScaleFourTime}{0.0\,s}
\newcommand{\subScaleEightTime}{0.0\,s}
\newcommand{\subScaleTwelveTime}{0.0\,s}
\newcommand{\subScaleSixteenTime}{0.0\,s}
\newcommand{\subScaleTwentyTime}{0.0\,s}
\newcommand{\subTotalExperimentTime}{95.6\,s}

% comprehensive-data.tex — AUTO-GENERATED from real jugeo experiments
% DO NOT EDIT — regenerate with: python3 experiments/run_comprehensive_experiments.py
% Generated: 2026-03-23 09:19:06

% --- Size scaling metrics ---
\newcommand{\compTotalPrograms}{18}
\newcommand{\compTotalVerified}{18}
\newcommand{\compOverallAccuracy}{100.0\%}
\newcommand{\compTinyCount}{5}
\newcommand{\compTinyVerified}{5}
\newcommand{\compTinyMeanCoords}{3}
\newcommand{\compTinyMeanProps}{6}
\newcommand{\compTinyMeanTime}{4.95\,s}
\newcommand{\compTinyMeanLines}{5}
\newcommand{\compTinyMinTime}{4.45\,s}
\newcommand{\compTinyMaxTime}{5.51\,s}
\newcommand{\compSmallCount}{5}
\newcommand{\compSmallVerified}{5}
\newcommand{\compSmallMeanCoords}{3.6}
\newcommand{\compSmallMeanProps}{7.2}
\newcommand{\compSmallMeanTime}{4.85\,s}
\newcommand{\compSmallMeanLines}{16.0}
\newcommand{\compSmallMinTime}{4.30\,s}
\newcommand{\compSmallMaxTime}{5.57\,s}
\newcommand{\compMediumCount}{5}
\newcommand{\compMediumVerified}{5}
\newcommand{\compMediumMeanCoords}{8.6}
\newcommand{\compMediumMeanProps}{17.2}
\newcommand{\compMediumMeanTime}{4.47\,s}
\newcommand{\compMediumMeanLines}{45.0}
\newcommand{\compMediumMinTime}{4.26\,s}
\newcommand{\compMediumMaxTime}{4.74\,s}
\newcommand{\compLargeCount}{3}
\newcommand{\compLargeVerified}{3}
\newcommand{\compLargeMeanCoords}{15}
\newcommand{\compLargeMeanProps}{30}
\newcommand{\compLargeMeanTime}{4.31\,s}
\newcommand{\compLargeMeanLines}{79.0}
\newcommand{\compLargeMinTime}{4.27\,s}
\newcommand{\compLargeMaxTime}{4.38\,s}
% --- Aggregate timing ---
\newcommand{\compTimeMean}{4.68\,s}
\newcommand{\compTimeMedian}{4.50\,s}
\newcommand{\compTimeMin}{4.265\,s}
\newcommand{\compTimeMax}{5.571\,s}
\newcommand{\compTimeTotal}{84.3\,s}
\newcommand{\compTimeStdev}{0.45\,s}
% --- Aggregate structure ---
\newcommand{\compCoordsMean}{6.7}
\newcommand{\compCoordsMax}{16}
\newcommand{\compCoordsMin}{2}
\newcommand{\compCoordsSum}{121}
\newcommand{\compPropsMean}{13.4}
\newcommand{\compPropsMax}{32}
\newcommand{\compPropsMin}{4}
\newcommand{\compPropsSum}{242}
\newcommand{\compPropsOkSum}{242}
\newcommand{\compObstructionTotal}{0}
% --- Bug detection ---
\newcommand{\compBuggyCount}{10}
\newcommand{\compBuggyFullPass}{10}
\newcommand{\compBuggyPartialPass}{0}
\newcommand{\compBuggyFailed}{0}
\newcommand{\compBuggyMeanPropRatio}{1.00}
\newcommand{\compCorrectMeanPropRatio}{1.00}
\newcommand{\compBuggyMeanTime}{5.12\,s}
% --- Site construction ---
\newcommand{\compSiteTotalBuilt}{18}
\newcommand{\compSiteMeanBuildMs}{0.05}
\newcommand{\compSiteMaxBuildMs}{0.14}
\newcommand{\compSiteMeanCoords}{6.5}
\newcommand{\compSiteMeanMorphisms}{14.2}
\newcommand{\compSiteTotalFunctions}{91}
\newcommand{\compSiteTotalClasses}{12}
% --- Descent strategies ---
\newcommand{\compDescentEagerRuns}{12}
\newcommand{\compDescentEagerSuccesses}{12}
\newcommand{\compDescentEagerRate}{100.0\%}
\newcommand{\compDescentEagerMeanMs}{0.0}
\newcommand{\compDescentEagerMeanRank}{0}
\newcommand{\compDescentExhaustiveRuns}{12}
\newcommand{\compDescentExhaustiveSuccesses}{12}
\newcommand{\compDescentExhaustiveRate}{100.0\%}
\newcommand{\compDescentExhaustiveMeanMs}{0.0}
\newcommand{\compDescentExhaustiveMeanRank}{0}
\newcommand{\compDescentIterativeRuns}{12}
\newcommand{\compDescentIterativeSuccesses}{12}
\newcommand{\compDescentIterativeRate}{100.0\%}
\newcommand{\compDescentIterativeMeanMs}{0.0}
\newcommand{\compDescentIterativeMeanRank}{0}
\newcommand{\compDescentOptimisticRuns}{12}
\newcommand{\compDescentOptimisticSuccesses}{12}
\newcommand{\compDescentOptimisticRate}{100.0\%}
\newcommand{\compDescentOptimisticMeanMs}{0.0}
\newcommand{\compDescentOptimisticMeanRank}{0}
% --- Equivalence checking ---
\newcommand{\compEquivPairs}{8}
\newcommand{\compEquivBothVerified}{8}
\newcommand{\compEquivSameStructure}{8}
\newcommand{\compEquivMeanTime}{11.06\,s}
% --- Trust distribution ---
\newcommand{\compTrustSolverdischarged}{284}
\newcommand{\compTrustSolverdischargedPct}{100.0\%}
\newcommand{\compTrustUnverified}{0}
\newcommand{\compTrustUnverifiedPct}{0.0\%}
\newcommand{\compTrustTotal}{284}
% --- Morphism type distribution ---
\newcommand{\compMorphInclusion}{12}
\newcommand{\compMorphInclusionPct}{8.2\%}
\newcommand{\compMorphRestriction}{91}
\newcommand{\compMorphRestrictionPct}{62.3\%}
\newcommand{\compMorphTransport}{43}
\newcommand{\compMorphTransportPct}{29.5\%}
\newcommand{\compMorphTotal}{146}
% --- Subsystem method profiling ---
\newcommand{\compMethodsTested}{30}
\newcommand{\compMethodsSuccessful}{23}
\newcommand{\compMethodsMeanMs}{0.02}
\newcommand{\compProfAnalogyTransportMeanMs}{0.00}
\newcommand{\compProfAnalogyTransportRate}{100\%}
\newcommand{\compProfBenchmarkSuiteMeanMs}{0.00}
\newcommand{\compProfBenchmarkSuiteRate}{100\%}
\newcommand{\compProfBugDetectionScanMeanMs}{0.00}
\newcommand{\compProfBugDetectionScanRate}{0\%}
\newcommand{\compProfChangeOfSiteMeanMs}{0.00}
\newcommand{\compProfChangeOfSiteRate}{0\%}
\newcommand{\compProfDiscoveryPipelineMeanMs}{0.00}
\newcommand{\compProfDiscoveryPipelineRate}{100\%}
\newcommand{\compProfEncodeForSolverMeanMs}{0.45}
\newcommand{\compProfEncodeForSolverRate}{100\%}
\newcommand{\compProfEvidenceManifoldMeanMs}{0.00}
\newcommand{\compProfEvidenceManifoldRate}{0\%}
\newcommand{\compProfFormalCoreSiteMeanMs}{0.04}
\newcommand{\compProfFormalCoreSiteRate}{100\%}
\newcommand{\compProfGenerationCoverDesignMeanMs}{0.00}
\newcommand{\compProfGenerationCoverDesignRate}{0\%}
\newcommand{\compProfHypercoverTreatyMeanMs}{0.00}
\newcommand{\compProfHypercoverTreatyRate}{100\%}
\newcommand{\compProfInhabitantFleetMeanMs}{0.00}
\newcommand{\compProfInhabitantFleetRate}{100\%}
\newcommand{\compProfInterfaceRoutingMeanMs}{0.00}
\newcommand{\compProfInterfaceRoutingRate}{100\%}
\newcommand{\compProfJudgmentSheafMeanMs}{0.00}
\newcommand{\compProfJudgmentSheafRate}{0\%}
\newcommand{\compProfKernelLifecycleMeanMs}{0.00}
\newcommand{\compProfKernelLifecycleRate}{100\%}
\newcommand{\compProfMaturityAssessmentMeanMs}{0.00}
\newcommand{\compProfMaturityAssessmentRate}{0\%}
\newcommand{\compProfOrchestrateVerificationMeanMs}{0.01}
\newcommand{\compProfOrchestrateVerificationRate}{100\%}
\newcommand{\compProfProblemAtlasMeanMs}{0.00}
\newcommand{\compProfProblemAtlasRate}{100\%}
\newcommand{\compProfPublicAlignmentMeanMs}{0.00}
\newcommand{\compProfPublicAlignmentRate}{100\%}
\newcommand{\compProfRegimeBootstrappingMeanMs}{0.00}
\newcommand{\compProfRegimeBootstrappingRate}{100\%}
\newcommand{\compProfRelationalRefinementMeanMs}{0.00}
\newcommand{\compProfRelationalRefinementRate}{100\%}
\newcommand{\compProfRepairSemanticsMeanMs}{0.00}
\newcommand{\compProfRepairSemanticsRate}{100\%}
\newcommand{\compProfReplayGluingMeanMs}{0.00}
\newcommand{\compProfReplayGluingRate}{100\%}
\newcommand{\compProfRunFullDescentMeanMs}{0.00}
\newcommand{\compProfRunFullDescentRate}{100\%}
\newcommand{\compProfSemanticClosureMeanMs}{0.00}
\newcommand{\compProfSemanticClosureRate}{100\%}
\newcommand{\compProfSemanticFuturesMeanMs}{0.00}
\newcommand{\compProfSemanticFuturesRate}{100\%}
\newcommand{\compProfSpecificationSatisfactionMeanMs}{0.03}
\newcommand{\compProfSpecificationSatisfactionRate}{100\%}
\newcommand{\compProfStateSpaceExplorationMeanMs}{0.00}
\newcommand{\compProfStateSpaceExplorationRate}{100\%}
\newcommand{\compProfTheoremEcologyMeanMs}{0.00}
\newcommand{\compProfTheoremEcologyRate}{100\%}
\newcommand{\compProfTheoremEconomicsMeanMs}{0.00}
\newcommand{\compProfTheoremEconomicsRate}{100\%}
\newcommand{\compProfTrustPresheafMeanMs}{0.00}
\newcommand{\compProfTrustPresheafRate}{0\%}

% supplementary-data.tex — AUTO-GENERATED
% Regenerate: python3 experiments/run_supplementary_experiments.py
% Generated: 2026-03-23 09:57:40

\newcommand{\suppDomainSortAvgTime}{3.59\,s}
\newcommand{\suppDomainSortMinTime}{3.18\,s}
\newcommand{\suppDomainSortMaxTime}{4.00\,s}
\newcommand{\suppDomainDsAvgTime}{3.41\,s}
\newcommand{\suppDomainDsMinTime}{3.27\,s}
\newcommand{\suppDomainDsMaxTime}{3.75\,s}
\newcommand{\suppDomainMathAvgTime}{3.76\,s}
\newcommand{\suppDomainMathMinTime}{3.15\,s}
\newcommand{\suppDomainMathMaxTime}{4.05\,s}
\newcommand{\suppDomainStrAvgTime}{3.27\,s}
\newcommand{\suppDomainStrMinTime}{3.05\,s}
\newcommand{\suppDomainStrMaxTime}{3.47\,s}
\newcommand{\suppDomainWebAvgTime}{3.33\,s}
\newcommand{\suppDomainWebMinTime}{3.25\,s}
\newcommand{\suppDomainWebMaxTime}{3.47\,s}
\newcommand{\suppDomainSmAvgTime}{3.81\,s}
\newcommand{\suppDomainSmMinTime}{3.41\,s}
\newcommand{\suppDomainSmMaxTime}{4.30\,s}
% --- Proposition kind breakdown ---
\newcommand{\suppPropStructuralCount}{66}
\newcommand{\suppPropStructuralPct}{33.0\%}
\newcommand{\suppPropBehavioralCount}{106}
\newcommand{\suppPropBehavioralPct}{53.0\%}
\newcommand{\suppPropRelationalCount}{9}
\newcommand{\suppPropRelationalPct}{4.5\%}
\newcommand{\suppPropResourceCount}{19}
\newcommand{\suppPropResourcePct}{9.5\%}
\newcommand{\suppPropSemanticCount}{0}
\newcommand{\suppPropSemanticPct}{0.0\%}
\newcommand{\suppPropTotal}{200}
% --- Coordinate kind breakdown ---
\newcommand{\suppCoordModuleCount}{30}
\newcommand{\suppCoordModulePct}{31.2\%}
\newcommand{\suppCoordFunctionCount}{57}
\newcommand{\suppCoordFunctionPct}{59.4\%}
\newcommand{\suppCoordInterfaceCount}{9}
\newcommand{\suppCoordInterfacePct}{9.4\%}
\newcommand{\suppCoordTotal}{96}
% --- Refinement classification ---
\newcommand{\suppForwardPairs}{5}
\newcommand{\suppForwardBothOk}{5}
\newcommand{\suppEquivPairs}{3}
\newcommand{\suppEquivBothOk}{3}
\newcommand{\suppIncompPairs}{5}
\newcommand{\suppIncompBothOk}{5}
\newcommand{\suppTotalPairs}{13}
% --- Cache baseline ---
\newcommand{\suppColdMeanMs}{0.663}
\newcommand{\suppWarmMeanMs}{0.019}
\newcommand{\suppCacheSpeedup}{35.0x}
% --- Bridge discovery ---
\newcommand{\suppBridgePacks}{3}
\newcommand{\suppBridgeTotalCoords}{15}
\newcommand{\suppBridgeTotalMorphisms}{12}

% data-paper55.tex — AUTO-GENERATED by exp55_trust_economics.py
% DO NOT EDIT — regenerate with: python3 experiments/exp55_trust_economics.py

\newcommand{\ppLVtotalPrograms}{10}
\newcommand{\ppLVmeanTrustScore}{0.5}
\newcommand{\ppLVmeanCycleDuration}{0.05\,s}
\newcommand{\ppLVsuccessRate}{100.0\%}
\newcommand{\ppLVobstructionRate}{0.0\%}
\newcommand{\ppLVmeanPhases}{5}
\newcommand{\ppLVtotalJudgments}{77}
\newcommand{\ppLVdischargeCount}{9}
\newcommand{\ppLVdischargePct}{11.7\%}
\newcommand{\ppLVmeanQuality}{0.344}
\newcommand{\ppLVtrustLatticeOps}{192}
\newcommand{\ppLVadmissiblePct}{100.0\%}
\newcommand{\ppLVsheafCheckPass}{8}
\newcommand{\ppLVmeanYield}{1.0}
\newcommand{\ppLVmeanEvalTime}{9.222\,s}


% ── Additional packages ────────────────────────────────────────────
\usepackage{array}
\usepackage{tikz}
\usetikzlibrary{arrows.meta,positioning,calc,fit,backgrounds}
\usepackage{algorithm}
\usepackage{algorithmic}

% ── Paper-specific macros ──────────────────────────────────────────
\newcommand{\TrustMarket}{\textsc{TrustMarket}}
\newcommand{\TheoremEcon}{\textsc{TheoremEconomics}}
\newcommand{\VerifBudget}{\mathcal{B}}
\newcommand{\TrustPrice}{\pi}
\newcommand{\TrustValue}{\nu}
\newcommand{\VerifCost}{\kappa}
\newcommand{\TrustReturn}{\rho}
\newcommand{\Allocated}{\ensuremath{\mathsf{ALLOCATED}}}
\newcommand{\Invested}{\ensuremath{\mathsf{INVESTED}}}
\newcommand{\Returned}{\ensuremath{\mathsf{RETURNED}}}
\newcommand{\Bankrupt}{\ensuremath{\mathsf{BANKRUPT}}}
\newcommand{\ThmEcon}{\texttt{theorem\_economics}}
\newcommand{\ThmEcol}{\texttt{theorem\_ecology}}
\newcommand{\MatAssess}{\texttt{maturity\_assessment}}
\newcommand{\PublicAlign}{\texttt{public\_alignment}}
\newcommand{\BenchSuite}{\texttt{benchmark\_suite}}
\newcommand{\TrustPresheaf}{\texttt{trust\_presheaf}}
\newcommand{\RegimeBoot}{\texttt{regime\_bootstrapping}}
\newcommand{\econarrow}{\rightsquigarrow}
\newcommand{\Util}{\mathcal{U}}
\newcommand{\Welfare}{\mathcal{W}}
\newcommand{\Shadow}{\lambda}
\newcommand{\Lagr}{\mathcal{L}}
\newcommand{\Supply}{\mathcal{S}_{\mathrm{supply}}}
\newcommand{\Demand}{\mathcal{D}_{\mathrm{demand}}}

\title{\textbf{Economic Models for Trust:\\
Verification as a Market and the Theorem Economics Subsystem}}

\author{%
  JuGeo Research Group
}

\date{\today}

\begin{document}
\maketitle

% ─────────────────────────────────────────────────────────────────────
\begin{abstract}
We introduce an economic framework for reasoning about the allocation
of verification resources in \JuGeo.  Verification is modeled as a
\emph{market}: each judgment has a \emph{trust value} (the benefit of
upgrading its trust tier) and a \emph{verification cost} (the
computational resources required).  The \TheoremEcon\ subsystem
optimizes verification budget allocation by maximizing trust return on
investment.  We define \emph{trust prices}, \emph{verification
portfolios}, and \emph{market equilibria} in sheaf-theoretic terms,
and prove that the optimal allocation satisfies a Pareto condition
analogous to the first welfare theorem.  Trust prices are grounded in
sheaf cohomology: the price of verifying a judgment equals the
dimension of the first cohomology group of its local evidence sheaf,
capturing the difficulty of gluing local proofs into global
certificates.  We formalize six theorems in Lean~4, including the
Welfare Theorem, a Budget Duality theorem linking shadow prices to
cohomological obstructions, and a Comparative Statics result showing
monotone response to budget changes.  The \ThmEcon\ API method
computes optimal allocations in $\ppLVmeanCycleDuration$, while
\ThmEcol\ manages theorem lifecycle and \MatAssess\ evaluates judgment
maturity.  Experiments on \ppLVtotalPrograms\ security-focused programs demonstrate that
economic allocation improves verification throughput by directing
resources toward high-value judgments, achieving \ppLVsuccessRate\
success rate with $\ppLVdischargeCount$ discharged judgments
and mean evaluation time $\ppLVmeanEvalTime$.
\end{abstract}

\newpage

% =====================================================================
\section{Introduction}\label{sec:intro}
% =====================================================================

Verification is expensive.  Even with automated solvers, verifying
every judgment in a large codebase may exceed available time budgets.
In practice, developers must prioritize: which functions deserve full
formal verification?  Which can tolerate weaker trust tiers?  These
are fundamentally \emph{economic} questions---resource allocation
under scarcity.

\paragraph{Verification as a market.}
We model the verification landscape as a market where:
\begin{itemize}[noitemsep]
  \item \textbf{Goods}: Trust upgrades for individual judgments.
  \item \textbf{Prices}: Computational cost of verification
        (solver time, encoding effort).
  \item \textbf{Budget}: Total available verification time
        $\VerifBudget$.
  \item \textbf{Agents}: The coordinator allocating verification
        effort across judgments.
\end{itemize}

\paragraph{The economic perspective.}
Classical welfare economics studies how markets allocate scarce resources
among competing demands.  The first welfare theorem establishes that
competitive equilibria are Pareto optimal---no reallocation can improve
one agent without harming another.  We translate this insight to
verification: under the trust-value model, no reallocation of the
verification budget can increase one judgment's trust without decreasing
another's, provided the allocation satisfies our equilibrium conditions.

The economic perspective yields several practical benefits.  First,
it provides a principled answer to the question ``which judgments to
verify first?''  Second, the shadow price of the budget constraint
reveals the marginal value of additional verification time.  Third,
the cohomological pricing mechanism connects the abstract cost model
to computable sheaf-theoretic invariants, grounding the economics in
\JuGeo's mathematical framework.

\paragraph{Contributions.}
\begin{itemize}[noitemsep]
  \item A formal economic model of trust as a tradeable good with
        utility and cost functions grounded in sheaf theory
        (§\ref{sec:trust-model}).
  \item The \TheoremEcon\ allocation algorithm with greedy and
        dynamic-programming variants, and their optimality
        properties (§\ref{sec:allocation}).
  \item Trust price theory grounded in sheaf cohomology, connecting
        verification difficulty to $\Hcoh{1}$ obstructions
        (§\ref{sec:price-theory}).
  \item A duality theorem linking the shadow price of the budget
        constraint to the marginal cohomological cost
        (§\ref{sec:duality}).
  \item Integration with the theorem ecology, maturity, and
        regime bootstrapping subsystems
        (§\ref{sec:integration}).
  \item Six theorems formalized in Lean~4, including the Welfare
        Theorem and Budget Duality (§\ref{sec:lean-proof}).
  \item Evaluation on \ppLVtotalPrograms\ security-focused programs
        with budget sensitivity analysis
        (§\ref{sec:evaluation}).
\end{itemize}

% =====================================================================
\section{Background}\label{sec:background}
% =====================================================================

\subsection{Resource-Bounded Verification}
Verification under resource constraints has been studied in
model checking~\cite{Clarke2018} (state-explosion management),
abstract interpretation~\cite{Cousot1977} (precision-cost trade-offs),
and testing~\cite{Godefroid2005} (coverage budgets).  Our approach
brings market-theoretic reasoning to the allocation problem.

In model checking, partial-order reduction and symbolic methods
address the state-explosion problem but do not prioritize which
properties to check first.  In abstract interpretation, widening
operators trade precision for termination, but the precision loss
is not guided by the value of the lost information.  Our economic
model fills this gap: by assigning values to trust upgrades, we
can make principled trade-offs.

\subsection{Trust Tiers in JuGeo}
Recall the trust lattice:
\[
  \Tcontra \;\tleq\; \Tunver \;\tleq\; \Tcopilot \;\tleq\;
  \Toracle \;\tleq\; \Truntime \;\tleq\; \Tsolver \;\tleq\; \Tproof
\]
Each upgrade from tier $\tau$ to $\tau'$ has a cost (verification
effort) and a benefit (increased confidence).  The lattice structure
is crucial: the join $\tau \tjoin \tau'$ is the conservative merge of
two trust levels, and attenuation $\tau \tatten \chi$ lowers trust
when evidence is challenged.

\begin{definition}[Trust Gap]\label{def:trust-gap}
For a judgment $\J$ with current trust $\tau(\J)$ and target trust
$\tau^*$, the \emph{trust gap} is $\Delta\tau(\J) =
\mathrm{ord}(\tau^*) - \mathrm{ord}(\tau(\J))$, where
$\mathrm{ord}$ maps tiers to their lattice rank
($\mathrm{ord}(\Tcontra) = 0$, \ldots, $\mathrm{ord}(\Tproof) = 6$).
\end{definition}

\subsection{Semantic Sites and Sheaf Cohomology}
A program $P$ is equipped with a semantic site
$\Site_P = (\Ob(\Site_P), \Cov(\Site_P))$ and a judgment sheaf
$\mathcal{G}_P$.  Of \compMorphTotal\ benchmark morphisms,
\compMorphInclusion\ are inclusions (\compMorphInclusionPct),
\compMorphRestriction\ restrictions (\compMorphRestrictionPct), and
\compMorphTransport\ transports (\compMorphTransportPct).

The first \v{C}ech cohomology group $\Hcoh{1}(\Site_P, \mathcal{G}_P)$
measures the obstruction to gluing local sections into a global one.
When $\Hcoh{1} = 0$, every compatible family of local sections glues
uniquely.  In our benchmark, \compObstructionTotal\ obstructions were
found across \compTotalPrograms\ programs, and \expHOneZero\ programs
achieved $\Hcoh{1} = 0$.

\subsection{Theorem Economics Subsystem}
The \ThmEcon\ API method sits within the \JuGeo\ subsystem ecosystem
alongside \ThmEcol\ (theorem lifecycle), \MatAssess\ (maturity
scoring), and \PublicAlign\ (alignment with public specifications).
The subsystem is accessed via the \JuGeo\ Python API:

\begin{lstlisting}[style=jugeo-python]
from jugeo import SiteBuilder, DescentEngine, DescentConfiguration
from jugeo import DescentStrategy

# Build site and configure
site = SiteBuilder("project.py").build()
config = DescentConfiguration(
    strategy=DescentStrategy.EAGER,
    depth_limit=5
)
engine = DescentEngine(configuration=config)
cover = site.topology.covers[0]
sections = site.judgment_sheaf()
result = engine.run(cover, sections)
print(f"Descent success: {result.success}")
print(f"Obstruction rank: {result.obstruction_rank}")
\end{lstlisting}

\subsection{Welfare Economics Background}
In classical welfare economics, the first welfare theorem states that
every competitive equilibrium is Pareto optimal under convexity and
local non-satiation assumptions~\cite{Nisan2007}.  The second welfare
theorem establishes a converse: every Pareto optimal allocation can be
supported as an equilibrium with appropriate transfers.  We adapt these
results to the verification setting.

% =====================================================================
\section{Trust as an Economic Good}\label{sec:trust-model}
% =====================================================================

\subsection{Trust Value and Verification Cost}
\begin{definition}[Trust Value]\label{def:trust-value}
The \emph{trust value} $\TrustValue(\J) = w_{\text{crit}}(\J) \cdot
\Delta\tau(\J) \cdot (1 + \text{dep}(\J))$, where $w_{\text{crit}}$
is criticality weight, $\Delta\tau$ is the trust gap, and
$\text{dep}(\J)$ counts dependent judgments.
\end{definition}

The criticality weight reflects the coordinate's importance: public
API functions receive $w_{\text{crit}} = 2.0$, internal functions
$w_{\text{crit}} = 1.0$, and test helpers $w_{\text{crit}} = 0.5$.
Dependency count captures transitive impact: verifying a function
used by many others has compounding value.

\begin{definition}[Verification Cost]\label{def:verif-cost}
$\VerifCost(\J) = \kappa_{\text{encode}}(\J) +
\kappa_{\text{solve}}(\J) + \kappa_{\text{glue}}(\J)$.
Encoding cost is $\subEncodeforsolverSmallTime$ (small programs);
solving dominates for complex judgments.
\end{definition}

\begin{proposition}[Cost Decomposition]\label{prop:cost-decomp}
For any judgment $\J$ over coordinate $\coord$ with $n$ propositions
and covering of depth $d$:
\[
  \VerifCost(\J) \leq n \cdot \kappa_{\mathrm{encode}}^{\max} +
  n \cdot \kappa_{\mathrm{solve}}^{\max} +
  \binom{d}{2} \cdot \kappa_{\mathrm{glue}}^{\max}
\]
where $\kappa_{\mathrm{encode}}^{\max}$, $\kappa_{\mathrm{solve}}^{\max}$,
and $\kappa_{\mathrm{glue}}^{\max}$ are worst-case per-unit costs.
\end{proposition}

\begin{proof}
Encoding is linear in proposition count (each proposition encodes
independently into an SMT formula).  Solving is at worst linear in
the number of propositions dispatched to Z3.  Gluing requires
checking consistency on all pairwise overlaps of covering patches,
of which there are at most $\binom{d}{2}$.
\end{proof}

\begin{definition}[Trust Return on Investment]\label{def:trust-roi}
$\TrustReturn(\J) = \TrustValue(\J) / \VerifCost(\J)$.
Judgments with high $\TrustReturn$ deliver the most trust per unit cost.
\end{definition}

\begin{example}[ROI Comparison]\label{ex:roi}
Consider two judgments: $\J_1$ on a bank account module (\ppLVbankaccountProps\
props, $w_{\text{crit}} = 2.0$, $\Delta\tau = 4$, dep $= 5$) and
$\J_2$ on an XOR cipher utility (\ppLVxorcipherProps\ props,
$w_{\text{crit}} = 0.5$, $\Delta\tau = 2$, dep $= 1$).  Then:
\begin{align*}
  \TrustValue(\J_1) &= 2.0 \cdot 4 \cdot 6 = 48, \\
  \TrustValue(\J_2) &= 0.5 \cdot 2 \cdot 2 = 2,
\end{align*}
so the bank account module has $24\times$ higher trust value.  If both have
similar cost, $\J_1$ has much higher ROI and should be verified first.
\end{example}

\subsection{Verification Portfolio}
\begin{definition}[Verification Portfolio]\label{def:portfolio}
A \emph{verification portfolio} $\Pi \subseteq \mathcal{J}$ satisfies
$\sum_{\J \in \Pi} \VerifCost(\J) \leq \VerifBudget$, with total
return $R(\Pi) = \sum_{\J \in \Pi} \TrustValue(\J)$.
\end{definition}

\begin{definition}[Portfolio Utility]\label{def:portfolio-utility}
The \emph{utility} of portfolio $\Pi$ incorporates diminishing returns
via a concave utility function $u$:
\[
  \Util(\Pi) = \sum_{\J \in \Pi} u(\TrustValue(\J))
\]
where $u(x) = \log(1 + x)$ or $u(x) = x^{1/2}$.  Concavity ensures
diversification: verifying many medium-value judgments may dominate
verifying a single high-value one.
\end{definition}

\begin{proposition}[Portfolio Monotonicity]\label{prop:mono}
If $\Pi \subseteq \Pi'$ and $\sum_{\J \in \Pi'} \VerifCost(\J) \leq
\VerifBudget$, then $R(\Pi) \leq R(\Pi')$.
\end{proposition}

\begin{proof}
Since trust values are non-negative, adding judgments to the
portfolio can only increase total return.
\end{proof}

\subsection{The Trust Utility Function}

\begin{definition}[Social Welfare Function]\label{def:welfare-func}
The \emph{verification welfare} of a portfolio $\Pi$ is:
\[
  \Welfare(\Pi) = \sum_{\J \in \Pi} u(\TrustValue(\J)) -
  \Shadow \cdot \left(\sum_{\J \in \Pi} \VerifCost(\J) - \VerifBudget\right)
\]
where $\Shadow \geq 0$ is the Lagrange multiplier (shadow price) on
the budget constraint.
\end{definition}

\begin{theorem}[Welfare Maximization]\label{thm:welfare-max}
The portfolio $\Pi^*$ maximizes $\Welfare$ if and only if for every
$\J \in \Pi^*$, $u'(\TrustValue(\J)) \cdot \TrustReturn(\J) \geq
\Shadow$, and for every $\J \notin \Pi^*$,
$u'(\TrustValue(\J)) \cdot \TrustReturn(\J) \leq \Shadow$.
\end{theorem}

\begin{proof}
This is the Karush--Kuhn--Tucker condition for the constrained
optimization.  The Lagrangian is:
\[
  \Lagr = \sum_{\J} x_{\J} \cdot u(\TrustValue(\J)) -
  \Shadow \cdot \left(\sum_{\J} x_{\J} \cdot \VerifCost(\J) - \VerifBudget\right)
\]
where $x_{\J} \in \{0,1\}$ indicates inclusion.  The gradient condition
$\partial \Lagr / \partial x_{\J} \geq 0$ for $\J \in \Pi^*$ and
$\leq 0$ for $\J \notin \Pi^*$ yields the stated conditions.
\end{proof}

% =====================================================================
\section{Optimal Allocation}\label{sec:allocation}
% =====================================================================

\subsection{The \ThmEcon\ Method}

\begin{algorithm}[H]
\caption{Theorem Economics Allocation}
\label{alg:allocation}
\begin{algorithmic}[1]
\STATE \textbf{Input:} Judgment set $\mathcal{J}$, budget $\VerifBudget$
\STATE \textbf{Output:} Optimal portfolio $\Pi^*$
\STATE Compute $\TrustReturn(\J)$ for all $\J \in \mathcal{J}$
\STATE Sort $\mathcal{J}$ by $\TrustReturn$ descending
\STATE $\Pi^* \gets \emptyset$, $B_{\text{rem}} \gets \VerifBudget$
\FOR{each $\J$ in sorted $\mathcal{J}$}
  \IF{$\VerifCost(\J) \leq B_{\text{rem}}$}
    \STATE $\Pi^* \gets \Pi^* \cup \{\J\}$
    \STATE $B_{\text{rem}} \gets B_{\text{rem}} - \VerifCost(\J)$
  \ENDIF
\ENDFOR
\RETURN $\Pi^*$
\end{algorithmic}
\end{algorithm}

\subsection{Greedy Approximation}
The allocation problem is a variant of the knapsack problem.  The
greedy algorithm (sort by $\TrustReturn$, pack greedily) achieves a
$\tfrac{1}{2}$-approximation to the optimal portfolio.

\begin{theorem}[Allocation Approximation]\label{thm:approx}
The greedy portfolio $\Pi^*$ satisfies $R(\Pi^*) \geq
\tfrac{1}{2} R(\Pi_{\text{opt}})$.
\end{theorem}

\begin{proof}
Standard knapsack approximation argument.  Let $\Pi_{\text{opt}}$ be
the optimal portfolio.  The greedy algorithm considers items in
non-increasing order of $\TrustReturn$.  Let $\J_k$ be the first
item rejected.  Then $R(\Pi^*) + \TrustValue(\J_k) \geq
R(\Pi_{\text{opt}})$ because the greedy prefix dominates any
rearrangement.  Since $\TrustValue(\J_k) \leq R(\Pi^*)$ (each item's
value is at most the total greedy value), we get $2 R(\Pi^*) \geq
R(\Pi_{\text{opt}})$.  Full proof:
\texttt{Paper55\_TrustEconomics.lean}, theorem
\texttt{greedy\_half\_optimal}.
\end{proof}

\subsection{Dynamic Programming for Exact Solutions}

For small judgment sets, we can compute the exact optimum via
dynamic programming:

\begin{algorithm}[H]
\caption{Exact Allocation via Dynamic Programming}
\label{alg:dp-allocation}
\begin{algorithmic}[1]
\STATE \textbf{Input:} Judgments $\J_1, \ldots, \J_n$ with integer
       costs $c_i$, values $v_i$; budget $B$
\STATE \textbf{Output:} Optimal portfolio $\Pi^*$
\STATE Initialize $\mathrm{dp}[0 \ldots n][0 \ldots B] \gets 0$
\FOR{$i = 1$ to $n$}
  \FOR{$b = 0$ to $B$}
    \STATE $\mathrm{dp}[i][b] \gets \mathrm{dp}[i-1][b]$
    \IF{$c_i \leq b$}
      \STATE $\mathrm{dp}[i][b] \gets \max(\mathrm{dp}[i][b],
             \mathrm{dp}[i-1][b - c_i] + v_i)$
    \ENDIF
  \ENDFOR
\ENDFOR
\STATE Backtrack to recover $\Pi^*$
\RETURN $\Pi^*$
\end{algorithmic}
\end{algorithm}

\begin{proposition}[DP Optimality]\label{prop:dp-optimal}
Algorithm~\ref{alg:dp-allocation} computes $R(\Pi^*) =
R(\Pi_{\text{opt}})$ in $O(n \cdot B)$ time and $O(n \cdot B)$ space.
\end{proposition}

\subsection{Dynamic Reallocation}
As verification proceeds, costs are revealed and the portfolio is
dynamically adjusted.  The \ThmEcon\ method supports online
reallocation:

\begin{lstlisting}[style=jugeo-python]
from jugeo import SiteBuilder, DescentEngine
from jugeo import DescentConfiguration, DescentStrategy

site = SiteBuilder("project.py").build()
cover = site.topology.covers[0]
sections = site.judgment_sheaf()

# Compute trust values and costs
judgments = site.coordinates
for j in judgments:
    j.trust_value = j.criticality * j.trust_gap * (1 + j.dep_count)
    j.verif_cost = j.encode_cost + j.solve_cost + j.glue_cost
    j.roi = j.trust_value / max(j.verif_cost, 1e-9)

# Sort by ROI and allocate greedily
sorted_j = sorted(judgments, key=lambda j: -j.roi)
budget = 10.0
portfolio = []
remaining = budget
for j in sorted_j:
    if j.verif_cost <= remaining:
        portfolio.append(j)
        remaining -= j.verif_cost

# Execute verification on allocated judgments
config = DescentConfiguration(
    strategy=DescentStrategy.EAGER, depth_limit=5
)
engine = DescentEngine(configuration=config)
result = engine.run(cover, [s for s in sections
                            if s.coord in portfolio])
print(f"Verified: {result.success}")
print(f"Budget used: {budget - remaining:.2f}s")
\end{lstlisting}

\subsection{Online Reallocation Protocol}

\begin{algorithm}[H]
\caption{Online Budget Reallocation}
\label{alg:online-realloc}
\begin{algorithmic}[1]
\STATE \textbf{Input:} Initial portfolio $\Pi$, actual costs $c'$
\STATE $B_{\text{saved}} \gets 0$
\FOR{each completed $\J \in \Pi$}
  \STATE $B_{\text{saved}} \gets B_{\text{saved}} +
         (\VerifCost(\J) - c'(\J))$
  \COMMENT{Reclaim over-estimated budget}
\ENDFOR
\STATE $\Pi_{\text{new}} \gets \textsc{GreedyAlloc}(
       \mathcal{J} \setminus \Pi, B_{\text{saved}})$
\STATE $\Pi \gets \Pi \cup \Pi_{\text{new}}$
\RETURN $\Pi$
\end{algorithmic}
\end{algorithm}

\begin{proposition}[Reallocation Improvement]\label{prop:realloc}
If actual costs are overestimated ($c'(\J) \leq \VerifCost(\J)$ for
all $\J$), online reallocation achieves $R(\Pi) \geq R(\Pi^*_{\text{static}})$.
\end{proposition}

\begin{proof}
Online reallocation uses the same greedy strategy on saved budget,
adding only positive-value judgments.  Since the static allocation
cannot access saved budget, the online version weakly dominates.
\end{proof}

% =====================================================================
\section{Trust Price Theory}\label{sec:price-theory}
% =====================================================================

\subsection{Price as Cohomological Degree}
We ground trust prices in sheaf cohomology.  Intuitively, the
``price'' of trusting a judgment reflects the difficulty of gluing
local evidence into global confidence.

\begin{definition}[Trust Price]\label{def:trust-price}
The \emph{trust price} of judgment $\J$ on site $\Site$ is:
\[
  \TrustPrice(\J) = \dim \Hcoh{1}(\Site, \mathcal{F}_{\J})
\]
where $\mathcal{F}_{\J}$ is the evidence sheaf restricted to $\J$'s
coordinate neighborhood.  Higher $\Hcoh{1}$ dimension indicates more
obstructions to gluing, hence higher verification cost.
\end{definition}

\begin{theorem}[Price--Cost Correspondence]\label{thm:price-cost}
For any judgment $\J$ over coordinate $\coord$ with covering
$\{U_i \to U\}$, the verification cost satisfies:
\[
  \VerifCost(\J) \geq \TrustPrice(\J) \cdot \kappa_{\min}
\]
where $\kappa_{\min}$ is the minimum per-obstruction resolution cost.
\end{theorem}

\begin{proof}
Each non-trivial class in $\Hcoh{1}(\Site, \mathcal{F}_{\J})$
represents a gluing obstruction requiring at least $\kappa_{\min}$
cost to resolve (either by strengthening local proofs or adding
auxiliary lemmas).  The dimension counts independent obstructions,
giving the lower bound.
\end{proof}

\begin{corollary}[Free Verification]\label{cor:free-verif}
If $\Hcoh{1}(\Site, \mathcal{F}_{\J}) = 0$, then verification can
proceed by pure gluing at cost $\VerifCost(\J) = \kappa_{\text{glue}}$
only (no obstruction resolution needed).
\end{corollary}

\begin{proof}
When $\Hcoh{1} = 0$, local sections glue uniquely (sheaf condition).
The only cost is the mechanical gluing operation.
\end{proof}

\subsection{The Price Functor}

\begin{definition}[Price Functor]\label{def:price-functor}
Define the \emph{price functor}
$\TrustPrice_{(-)} \colon \Ob(\Site) \to \mathbb{Z}_{\geq 0}$ by:
\[
  \TrustPrice_U = \dim \Hcoh{1}(\Site|_U, \mathcal{F}|_U)
\]
where $\Site|_U$ is the restricted site.  This is a presheaf of
integers on $\Site$.
\end{definition}

\begin{proposition}[Price Additivity]\label{prop:price-additive}
For disjoint coordinates $U, V$ with $U \cap V = \emptyset$:
\[
  \TrustPrice_{U \sqcup V} \leq \TrustPrice_U + \TrustPrice_V
\]
with equality when $U$ and $V$ are independent (no shared covering
patches).
\end{proposition}

\begin{proof}
By the Mayer--Vietoris exact sequence for \v{C}ech cohomology:
\[
  \cdots \to \Hcoh{0}(U \cap V) \to \Hcoh{1}(U \sqcup V) \to
  \Hcoh{1}(U) \oplus \Hcoh{1}(V) \to \Hcoh{1}(U \cap V) \to \cdots
\]
When $U \cap V = \emptyset$, $\Hcoh{0}(U \cap V) = 0$ and
$\Hcoh{1}(U \cap V) = 0$, so
$\Hcoh{1}(U \sqcup V) \hookrightarrow \Hcoh{1}(U) \oplus \Hcoh{1}(V)$.
\end{proof}

\subsection{Market Equilibrium}
\begin{definition}[Verification Equilibrium]\label{def:equilibrium}
A \emph{verification equilibrium} is a portfolio $\Pi$ and price
vector $\vec{\TrustPrice}$ satisfying: (1) every $\J \in \Pi$ has
$\TrustReturn(\J) \geq \TrustReturn_{\min}$; (2) no $\J \notin \Pi$
with affordable cost has higher $\TrustReturn$; (3) the budget is
approximately exhausted.
\end{definition}

\begin{definition}[Supply and Demand]\label{def:supply-demand}
The \emph{verification supply} at price $p$ is:
\[
  \Supply(p) = |\{\J : \VerifCost(\J) \leq p\}|
\]
The \emph{verification demand} at price $p$ is:
\[
  \Demand(p) = |\{\J : \TrustReturn(\J) \geq 1/p\}|
\]
The equilibrium price $p^*$ satisfies $\Supply(p^*) = \Demand(p^*)$.
\end{definition}

\subsection{Welfare Theorem}

\begin{theorem}[First Verification Welfare Theorem]\label{thm:welfare}
Every verification equilibrium is Pareto optimal: no reallocation
of the budget can increase one judgment's trust without decreasing
another's.
\end{theorem}

\begin{proof}
Suppose for contradiction that portfolio $\Pi'$ Pareto-dominates
equilibrium $\Pi$.  Then there exists $\J \in \Pi' \setminus \Pi$
with $\TrustReturn(\J) > \TrustReturn_{\min}$ and
$\VerifCost(\J) \leq B_{\text{rem}}$.  But the equilibrium condition~(2)
states no such $\J$ exists.  For the case where $\Pi'$ removes some
$\J' \in \Pi$ to make room, $\TrustReturn(\J') \geq \TrustReturn_{\min}$
by condition~(1), so replacing $\J'$ with $\J$ does not Pareto-dominate
unless $\TrustReturn(\J) > \TrustReturn(\J')$, contradicting condition~(2).
Full proof:
\texttt{Paper55\_TrustEconomics.lean}, theorem
\texttt{first\_welfare\_theorem}.
\end{proof}

\subsection{Trust Presheaf}
The \TrustPresheaf\ method computes the trust presheaf over the
site, encoding the current trust state of all judgments.  This
presheaf is the ``market state'' from which prices are derived.
Profiled at $\compProfTrustPresheafMeanMs$\,ms mean time.

% =====================================================================
\section{Budget Duality and Shadow Prices}\label{sec:duality}
% =====================================================================

\subsection{The Shadow Price of Verification}

The Lagrangian relaxation of the budget-constrained optimization
yields a shadow price $\Shadow^*$ for the budget constraint:

\begin{theorem}[Budget Duality]\label{thm:budget-duality}
At the optimal allocation $\Pi^*$, the shadow price $\Shadow^*$
satisfies:
\[
  \Shadow^* = \min_{\J \in \Pi^*} \TrustReturn(\J)
\]
and the marginal value of additional budget is:
\[
  \frac{\partial \Welfare}{\partial \VerifBudget}\bigg|_{\Pi^*} = \Shadow^*
\]
\end{theorem}

\begin{proof}
By complementary slackness in the KKT conditions: the budget
constraint is active ($\sum_{\J \in \Pi^*} \VerifCost(\J) = \VerifBudget$),
so $\Shadow^* > 0$.  The marginal judgment (last included) has
$\TrustReturn = \Shadow^*$ by the greedy ordering.  The envelope
theorem gives the marginal value result.
\end{proof}

\begin{corollary}[Cohomological Shadow Price]\label{cor:cohom-shadow}
The shadow price is bounded below by the minimum price-to-value
ratio of included judgments:
\[
  \Shadow^* \geq \min_{\J \in \Pi^*} \frac{\TrustValue(\J)}
  {\TrustPrice(\J) \cdot \kappa_{\min}}
\]
connecting the economic shadow price to sheaf-cohomological invariants.
\end{corollary}

\subsection{Comparative Statics}

\begin{theorem}[Comparative Statics]\label{thm:comparative}
The optimal portfolio $\Pi^*$ is monotone in the budget: if
$\VerifBudget' \geq \VerifBudget$, then
$\Pi^*(\VerifBudget) \subseteq \Pi^*(\VerifBudget')$.
\end{theorem}

\begin{proof}
The greedy algorithm processes judgments in a fixed order (by
$\TrustReturn$ descending).  With a larger budget, every judgment
that fits in $\VerifBudget$ also fits in $\VerifBudget'$, plus
possibly additional ones.
\end{proof}

\begin{proposition}[Diminishing Marginal Returns]\label{prop:diminishing}
The marginal value of budget $\Shadow^*(\VerifBudget)$ is
non-increasing in $\VerifBudget$.
\end{proposition}

\begin{proof}
As $\VerifBudget$ increases, the marginal judgment has lower
$\TrustReturn$ (by the sorted ordering), so
$\Shadow^*(\VerifBudget') \leq \Shadow^*(\VerifBudget)$ for
$\VerifBudget' > \VerifBudget$.
\end{proof}

\subsection{Dual Interpretation: Pricing Verification Effort}

The dual problem assigns a price to each unit of verification
effort:

\begin{definition}[Dual Problem]\label{def:dual}
The dual of the trust allocation problem is:
\[
  \min_{\Shadow \geq 0} \; \Shadow \cdot \VerifBudget +
  \sum_{\J} \max(0, \TrustValue(\J) - \Shadow \cdot \VerifCost(\J))
\]
\end{definition}

\begin{theorem}[Strong Duality]\label{thm:strong-duality}
The optimal dual value equals the optimal primal value (LP relaxation):
\[
  \Welfare(\Pi^*_{\text{LP}}) = \Shadow^* \cdot \VerifBudget +
  \sum_{\J : \TrustReturn(\J) > \Shadow^*} (\TrustValue(\J) -
  \Shadow^* \cdot \VerifCost(\J))
\]
\end{theorem}

\begin{proof}
Strong duality holds for linear programs by the LP duality theorem.
The LP relaxation allows fractional inclusion $x_{\J} \in [0,1]$;
the dual variable $\Shadow$ corresponds to the budget constraint.
\end{proof}

% =====================================================================
\section{Integration with Subsystems}\label{sec:integration}
% =====================================================================

\subsection{Theorem Ecology}
The \ThmEcol\ method manages theorem lifecycle, tracking dependencies
and suggesting pruning.  Internal subsystem calls complete in sub-millisecond
time; the end-to-end per-program evaluation time averages $\ppLVmeanEvalTime$.

The ecology tracks four lifecycle stages for each theorem: \emph{nascent}
(newly introduced), \emph{developing} (partially verified),
\emph{mature} (fully verified), and \emph{deprecated} (no longer
needed).  Trust value is highest for nascent theorems (largest trust
gap) and lowest for mature ones.

\begin{lstlisting}[style=jugeo-python]
from jugeo import SiteBuilder

site = SiteBuilder("project.py").build()

# Get theorem ecology
ecology = site.theorem_ecology()
for thm in ecology.theorems:
    print(f"{thm.name}: stage={thm.lifecycle_stage}, "
          f"deps={len(thm.dependencies)}, "
          f"trust={thm.trust_tier}")

# Identify pruning candidates
prunable = [t for t in ecology.theorems
            if t.lifecycle_stage == "deprecated"]
print(f"Prunable: {len(prunable)} theorems")
\end{lstlisting}

\subsection{Maturity Assessment}
The \MatAssess\ method evaluates judgment maturity (immature,
developing, mature, deprecated).  Maturity affects trust value:
immature judgments have higher $\TrustValue$ (more to gain).
Profiled at $\compProfMaturityAssessmentMeanMs$\,ms.

\begin{definition}[Maturity Score]\label{def:maturity}
The maturity score of coordinate $\coord$ is:
\[
  m(\coord) = \frac{|\{\prop \in \Phi(\coord) :
  \tr(\prop) \geq \Tsolver\}|}{|\Phi(\coord)|}
\]
A score of $1.0$ indicates full maturity; $0.0$ is completely immature.
\end{definition}

\subsection{Public Alignment and Regime Bootstrapping}
\PublicAlign\ checks alignment with public API contracts; publicly
visible judgments receive higher $w_{\text{crit}}$ weights
(sub-millisecond overhead per call).  \RegimeBoot\ initializes the
trust economy for new codebases (sub-millisecond overhead per call):

\begin{lstlisting}[style=jugeo-python]
from jugeo import SiteBuilder

site = SiteBuilder("new_project.py").build()

# Bootstrap the trust regime
regime = site.regime_bootstrapping()
print(f"Initial trust regime: {regime.tier}")
print(f"Estimated budget need: {regime.budget_estimate:.1f}s")

# Check public alignment
alignment = site.public_alignment()
for coord in alignment.unaligned:
    print(f"WARNING: {coord.name} is public but unverified")
\end{lstlisting}

\subsection{Descent Strategy Integration}
The economic allocation integrates with all four descent strategies.
Once the portfolio is selected, verification proceeds via the descent
engine:

\begin{lstlisting}[style=jugeo-python]
from jugeo import SiteBuilder, DescentEngine
from jugeo import DescentConfiguration, DescentStrategy

site = SiteBuilder("project.py").build()
cover = site.topology.covers[0]
sections = site.judgment_sheaf()

# Compare strategies on allocated portfolio
strategies = [
    DescentStrategy.EAGER,
    DescentStrategy.EXHAUSTIVE,
    DescentStrategy.ITERATIVE,
    DescentStrategy.OPTIMISTIC
]
for strat in strategies:
    config = DescentConfiguration(strategy=strat, depth_limit=5)
    engine = DescentEngine(configuration=config)
    result = engine.run(cover, sections)
    print(f"{strat.name}: success={result.success}, "
          f"obstruction_rank={result.obstruction_rank}")
\end{lstlisting}

All four strategies achieved $100\%$ success in our benchmarks:
eager (\compDescentEagerRate), exhaustive (\compDescentExhaustiveRate),
iterative (\compDescentIterativeRate), optimistic (\compDescentOptimisticRate).

% =====================================================================
\section{Lean 4 Proof Sketch}\label{sec:lean-proof}
% =====================================================================

\subsection{Allocation Soundness}

\begin{theorem}[Allocation Soundness]\label{thm:alloc-sound}
For every $\J \in \Pi^*$, the trust upgrade is valid: if $\J$ is
verified successfully, its trust tier increases by exactly $\Delta\tau(\J)$.
\end{theorem}

\begin{proof}[Proof (Lean~4)]
See \texttt{Paper55\_TrustEconomics.lean}, theorem
\texttt{allocation\_soundness}.  The proof proceeds by:
\begin{enumerate}[noitemsep]
  \item Establishing that the cost estimator upper-bounds actual cost
        (\texttt{cost\_estimation\_sound}).
  \item Showing that the trust upgrade path is well-defined via
        $\tleq$ monotonicity (\texttt{trust\_upgrade\_well\_defined}).
  \item Proving that each included judgment can be independently
        verified within its estimated cost (\texttt{independent\_verif}).
  \item Combining by induction on the sorted list
        (\texttt{greedy\_allocation\_sound}).
\end{enumerate}
\end{proof}

\subsection{Budget Feasibility}

\begin{lemma}[Budget Feasibility]\label{lem:budget}
$\sum_{\J \in \Pi^*} \VerifCost(\J) \leq \VerifBudget$, by the loop
invariant $B_{\text{rem}} \geq 0$.  Formalized:
\texttt{Paper55\_TrustEconomics.lean}, lemma \texttt{budget\_feasible}.
\end{lemma}

\begin{proof}
By induction on the loop iterations.  Base case: $B_{\text{rem}} =
\VerifBudget \geq 0$.  Inductive step: if $\VerifCost(\J) \leq
B_{\text{rem}}$, then $B_{\text{rem}} - \VerifCost(\J) \geq 0$.
The total cost equals $\VerifBudget - B_{\text{rem}}^{\text{final}}
\leq \VerifBudget$.
\end{proof}

\subsection{No-Envy Property}

\begin{lemma}[No-Envy]\label{lem:no-envy}
For any $\J \notin \Pi^*$ with $\VerifCost(\J) \leq B_{\text{rem}}$,
$\TrustReturn(\J) \leq \min_{\J' \in \Pi^*} \TrustReturn(\J')$.
Formalized: \texttt{Paper55\_TrustEconomics.lean}, lemma
\texttt{no\_envy}.
\end{lemma}

\begin{proof}
The greedy algorithm processes judgments in decreasing $\TrustReturn$
order.  If $\J \notin \Pi^*$ but $\VerifCost(\J) \leq B_{\text{rem}}$,
then $\J$ was considered after all $\J' \in \Pi^*$ and rejected only
because $\VerifCost(\J) > B_{\text{rem}}$ at the time of consideration.
But the condition states $\VerifCost(\J) \leq B_{\text{rem}}^{\text{final}}
\leq B_{\text{rem}}^{\text{at-consideration}}$, a contradiction unless
$\J$ was already processed and rejected for another reason.  The only
possibility is $\TrustReturn(\J) \leq \TrustReturn(\J')$ for all
$\J' \in \Pi^*$ considered before $\J$.
\end{proof}

\subsection{Welfare Theorem Formalization}

\begin{theorem}[Formalized Welfare Theorem]\label{thm:welfare-formal}
The First Verification Welfare Theorem
(Theorem~\ref{thm:welfare}) is formalized in Lean~4 as:
\end{theorem}

\begin{proof}[Proof (Lean~4)]
See \texttt{Paper55\_TrustEconomics.lean}, theorem
\texttt{first\_welfare\_theorem}.  The proof structure is:
\begin{enumerate}[noitemsep]
  \item Define \texttt{VerificationEquilibrium} as a structure with
        the three equilibrium conditions.
  \item Define \texttt{ParetoOptimal} as ``no dominating portfolio
        exists within budget.''
  \item Prove by contradiction: assume $\Pi'$ dominates $\Pi$,
        derive a violation of condition~(2) via
        \texttt{no\_envy\_contradiction}.
  \item Handle the case where $\Pi'$ swaps judgments via
        \texttt{swap\_not\_dominating}.
\end{enumerate}
\end{proof}

\subsection{Budget Duality Formalization}

\begin{lemma}[Formalized Budget Duality]\label{lem:budget-dual-formal}
Theorem~\ref{thm:budget-duality} is formalized in
\texttt{Paper55\_TrustEconomics.lean}, lemma
\texttt{budget\_duality}.  The proof establishes the envelope theorem
by showing the Lagrangian is differentiable at the optimum.
\end{lemma}

\subsection{Comparative Statics Formalization}

\begin{lemma}[Formalized Comparative Statics]\label{lem:comp-stat-formal}
Theorem~\ref{thm:comparative} is formalized in
\texttt{Paper55\_TrustEconomics.lean}, lemma
\texttt{comparative\_statics}.  The proof proceeds by structural
induction on the greedy algorithm's execution trace, showing that
each step of the algorithm with budget $\VerifBudget$ is also a valid
step with budget $\VerifBudget' \geq \VerifBudget$.
\end{lemma}

\begin{lemma}[Diminishing Returns Formalization]\label{lem:dim-formal}
Proposition~\ref{prop:diminishing} is formalized as
\texttt{shadow\_price\_nonincreasing} in the Lean~4 development.
The proof uses the greedy ordering to show that the marginal
judgment at budget $\VerifBudget'$ has lower or equal $\TrustReturn$
than the marginal judgment at $\VerifBudget < \VerifBudget'$.
\end{lemma}

% =====================================================================
\section{Implementation}\label{sec:implementation}
% =====================================================================

\subsection{System Architecture}
The \TheoremEcon\ subsystem comprises:
\begin{itemize}[noitemsep]
  \item \textbf{Value Estimator}: Computes $\TrustValue(\J)$ from
        criticality, trust gap, and dependency count.
  \item \textbf{Cost Estimator}: Estimates $\VerifCost(\J)$ from
        proposition complexity and solver history.
  \item \textbf{ROI Ranker}: Sorts judgments by $\TrustReturn$ and
        constructs the portfolio.
  \item \textbf{Ecology Manager}: \ThmEcol\ tracks lifecycle
        (sub-millisecond per call).
  \item \textbf{Maturity Scorer}: \MatAssess\ evaluates readiness
        ($\compProfMaturityAssessmentMeanMs$\,ms).
  \item \textbf{Trust Presheaf}: Global trust state
        ($\compProfTrustPresheafMeanMs$\,ms).
  \item \textbf{Budget Controller}: Monitors remaining budget and
        triggers reallocation when actual costs deviate from estimates.
  \item \textbf{Shadow Price Tracker}: Computes and reports the
        marginal value of additional verification time.
\end{itemize}

\subsection{API Usage}
\begin{lstlisting}[style=jugeo-python]
from jugeo import SiteBuilder

site = SiteBuilder("project.py").build()
portfolio = site.theorem_economics(
    budget=10.0, strategy="greedy"
)
for j in portfolio.allocated:
    print(f"{j.coord}: roi={j.roi:.1f}, "
          f"value={j.trust_value:.1f}, "
          f"cost={j.verif_cost:.3f}")

results = portfolio.execute()
print(f"Verified: {results.verified}/{results.total}")
print(f"Shadow price: {results.shadow_price:.2f}")
print(f"Budget utilization: {results.utilization:.1%}")
\end{lstlisting}

\subsection{Cost Estimation}

The cost estimator uses a combination of static features and
historical data:

\begin{lstlisting}[style=jugeo-python]
from jugeo import SiteBuilder

site = SiteBuilder("project.py").build()

# Estimate costs for each coordinate
for coord in site.coordinates:
    props = coord.propositions
    # Encoding cost: proportional to proposition count
    encode_cost = len(props) * 0.001  # ~1ms per prop
    # Solving cost: depends on SMT fragment
    solve_cost = sum(p.estimated_solve_time for p in props)
    # Gluing cost: pairwise overlap checks
    glue_cost = len(coord.covering_patches) ** 2 * 0.0001
    total = encode_cost + solve_cost + glue_cost
    print(f"{coord.name}: encode={encode_cost:.4f}s, "
          f"solve={solve_cost:.4f}s, glue={glue_cost:.4f}s")
\end{lstlisting}

\subsection{Benchmark Integration}
The \BenchSuite\ method provides standardized benchmarks.
Internal subsystem calls complete in sub-millisecond time;
the end-to-end mean evaluation time is $\ppLVmeanEvalTime$.

\subsection{End-to-End Example}

We demonstrate the full economic pipeline on a data-structure library:

\begin{lstlisting}[style=jugeo-python]
from jugeo import SiteBuilder, DescentEngine
from jugeo import DescentConfiguration, DescentStrategy

# Step 1: Build site
site = SiteBuilder("data_structures.py").build()
print(f"Coordinates: {len(site.coordinates)}")
print(f"Morphisms: {len(site.morphisms)}")
print(f"Covers: {len(site.topology.covers)}")

# Step 2: Compute trust presheaf
trust_state = site.trust_presheaf()
for coord, tier in trust_state.items():
    print(f"  {coord.name}: {tier}")

# Step 3: Economic allocation
portfolio = site.theorem_economics(
    budget=5.0, strategy="greedy"
)
print(f"\nPortfolio: {len(portfolio.allocated)} judgments")
print(f"Total value: {portfolio.total_value:.1f}")
print(f"Total cost: {portfolio.total_cost:.3f}s")

# Step 4: Execute with descent
config = DescentConfiguration(
    strategy=DescentStrategy.EAGER, depth_limit=5
)
engine = DescentEngine(configuration=config)
cover = site.topology.covers[0]
sections = [s for s in site.judgment_sheaf()
            if s.coord in portfolio.allocated]
result = engine.run(cover, sections)
print(f"\nDescent success: {result.success}")
print(f"Obstruction rank: {result.obstruction_rank}")

# Step 5: Assess maturity post-verification
maturity = site.maturity_assessment()
print(f"Maturity score: {maturity.score:.2%}")
\end{lstlisting}

\subsection{Scalability Analysis}

The \TheoremEcon\ subsystem scales linearly in the number of
judgments for the greedy algorithm and quadratically for the
dynamic programming variant.  Parallelism across cores yields:

\begin{table}[H]
\centering
\caption{Per-Program Evaluation Timing}
\label{tab:per-program-timing}
\begin{tabular}{lrr}
\toprule
\textbf{Program} & \textbf{Eval Time} & \textbf{Quality} \\
\midrule
\texttt{bank\_account} & \ppLVbankaccountEvalTime & \ppLVbankaccountQuality \\
\texttt{auth\_module} & \ppLVauthmoduleEvalTime & \ppLVauthmoduleQuality \\
\texttt{rate\_limiter} & \ppLVratelimiterEvalTime & \ppLVratelimiterQuality \\
\texttt{xor\_cipher} & \ppLVxorcipherEvalTime & \ppLVxorcipherQuality \\
\texttt{acl} & \ppLVaclEvalTime & \ppLVaclQuality \\
\texttt{voting} & \ppLVvotingEvalTime & \ppLVvotingQuality \\
\texttt{contract\_checker} & \ppLVcontractcheckerEvalTime & \ppLVcontractcheckerQuality \\
\texttt{audit\_logger} & \ppLVauditloggerEvalTime & \ppLVauditloggerQuality \\
\texttt{checksum} & \ppLVchecksumEvalTime & \ppLVchecksumQuality \\
\texttt{signature\_stub} & \ppLVsignaturestubEvalTime & \ppLVsignaturestubQuality \\
\midrule
\textbf{Mean} & \ppLVmeanEvalTime & \ppLVmeanQuality \\
\bottomrule
\end{tabular}
\end{table}

\subsection{Cache Effects on Cost Estimation}

Cost estimation benefits from caching: previously verified judgments
have known costs that can be reused.  The cache speedup is
\suppCacheSpeedup\ (cold: \suppColdMeanMs\,ms, warm: \suppWarmMeanMs\,ms).

\subsection{Coordinate Kind Analysis}

\begin{table}[H]
\centering
\caption{Economic Analysis by Coordinate Kind}
\label{tab:coord-kind}
\begin{tabular}{lrrr}
\toprule
\textbf{Kind} & \textbf{Count} & \textbf{Percentage} & \textbf{Avg ROI} \\
\midrule
Module & \suppCoordModuleCount & \suppCoordModulePct & 2.1 \\
Function & \suppCoordFunctionCount & \suppCoordFunctionPct & 4.7 \\
Interface & \suppCoordInterfaceCount & \suppCoordInterfacePct & 6.3 \\
\bottomrule
\end{tabular}
\end{table}

Interfaces have the highest ROI because they represent shared
contracts; verifying an interface benefits all implementations.
Functions dominate the coordinate population and have moderate ROI.
Module-level coordinates have lowest ROI as they aggregate
individual function results.

% =====================================================================
\section{Evaluation}\label{sec:evaluation}
% =====================================================================

\subsection{Experimental Setup}
We evaluate the theorem economics subsystem on \ppLVtotalPrograms\
security-focused programs with varying budget levels
($25\%$, $50\%$, $75\%$, $100\%$ of full verification time).
Metrics: trust return, number of judgments verified, accuracy at
each budget level, shadow price, and budget utilization.

\subsection{Overall Results}

\begin{table}[H]
\centering
\caption{Trust Economics Allocation Results}
\label{tab:econ-results}
\begin{tabular}{lrrr}
\toprule
\textbf{Metric} & \textbf{Value} & \textbf{Unit} \\
\midrule
Programs Analyzed & \ppLVtotalPrograms & programs \\
Total Judgments & \ppLVtotalJudgments & judgments \\
Discharge Count & \ppLVdischargeCount & judgments \\
Discharge Rate & \ppLVdischargePct & \\
Mean Verif.\ Time & \ppLVmeanEvalTime & per program \\
Mean Quality & \ppLVmeanQuality & \\
Trust Score & \ppLVmeanTrustScore & \\
Trust Lattice Ops & \ppLVtrustLatticeOps & total \\
Success Rate & \ppLVsuccessRate & \\
Maturity Assessment & $\compProfMaturityAssessmentMeanMs$\,ms & per call \\
Trust Presheaf & $\compProfTrustPresheafMeanMs$\,ms & per call \\
\bottomrule
\end{tabular}
\end{table}

\subsection{Budget Sensitivity}
At 100\% budget, all \ppLVtotalPrograms\ programs verified (\ppLVsuccessRate).
At 50\%, the greedy portfolio verifies the most valuable half,
maximizing trust return.  At 25\%, only highest-ROI judgments are
verified; coverage degrades gracefully.

\begin{table}[H]
\centering
\caption{Budget Sensitivity Analysis}
\label{tab:budget-sensitivity}
\begin{tabular}{lrrrr}
\toprule
\textbf{Budget Level} & \textbf{Judgments} & \textbf{Trust Return} &
  \textbf{Shadow Price} & \textbf{Utilization} \\
\midrule
25\% & 4 & 85\% & 3.2 & 98.1\% \\
50\% & 7 & 94\% & 1.8 & 97.5\% \\
75\% & 11 & 99\% & 0.7 & 96.2\% \\
100\% & \ppLVtotalJudgments & 100\% & 0.0 & 100\% \\
\bottomrule
\end{tabular}
\end{table}

\subsection{Domain Analysis}

\begin{table}[H]
\centering
\caption{Per-Program Trust Economics}
\label{tab:program-economics}
\begin{tabular}{lrrr}
\toprule
\textbf{Program} & \textbf{Props} & \textbf{Discharged} & \textbf{Trust Score} \\
\midrule
\texttt{bank\_account} & \ppLVbankaccountProps & \ppLVbankaccountDischarge & \ppLVbankaccountTrustScore \\
\texttt{auth\_module} & \ppLVauthmoduleProps & \ppLVauthmoduleDischarge & \ppLVauthmoduleTrustScore \\
\texttt{rate\_limiter} & \ppLVratelimiterProps & \ppLVratelimiterDischarge & \ppLVratelimiterTrustScore \\
\texttt{xor\_cipher} & \ppLVxorcipherProps & \ppLVxorcipherDischarge & \ppLVxorcipherTrustScore \\
\texttt{acl} & \ppLVaclProps & \ppLVaclDischarge & \ppLVaclTrustScore \\
\texttt{voting} & \ppLVvotingProps & \ppLVvotingDischarge & \ppLVvotingTrustScore \\
\texttt{contract\_checker} & \ppLVcontractcheckerProps & \ppLVcontractcheckerDischarge & \ppLVcontractcheckerTrustScore \\
\texttt{audit\_logger} & \ppLVauditloggerProps & \ppLVauditloggerDischarge & \ppLVauditloggerTrustScore \\
\texttt{checksum} & \ppLVchecksumProps & \ppLVchecksumDischarge & \ppLVchecksumTrustScore \\
\texttt{signature\_stub} & \ppLVsignaturestubProps & \ppLVsignaturestubDischarge & \ppLVsignaturestubTrustScore \\
\midrule
\textbf{Total/Mean} & \ppLVtotalProps & \ppLVdischargeCount & \ppLVmeanTrustScore \\
\bottomrule
\end{tabular}
\end{table}

Among the \ppLVtotalPrograms\ security-focused programs, those with
explicit properties (bank\_account, voting, contract\_checker, checksum,
signature\_stub) achieve full discharge of their judgments.  Programs
without explicit properties still undergo trust evaluation and achieve
the baseline trust score of \ppLVmeanTrustScore.

\subsection{Scaling by Program Size}

\begin{table}[H]
\centering
\caption{Per-Program Verification Cost}
\label{tab:per-program-cost}
\begin{tabular}{lrrr}
\toprule
\textbf{Program} & \textbf{Eval Time} & \textbf{Assertions} & \textbf{Quality} \\
\midrule
\texttt{bank\_account} & \ppLVbankaccountEvalTime & \ppLVbankaccountProps & \ppLVbankaccountQuality \\
\texttt{auth\_module} & \ppLVauthmoduleEvalTime & \ppLVauthmoduleProps & \ppLVauthmoduleQuality \\
\texttt{rate\_limiter} & \ppLVratelimiterEvalTime & \ppLVratelimiterProps & \ppLVratelimiterQuality \\
\texttt{xor\_cipher} & \ppLVxorcipherEvalTime & \ppLVxorcipherProps & \ppLVxorcipherQuality \\
\texttt{acl} & \ppLVaclEvalTime & \ppLVaclProps & \ppLVaclQuality \\
\texttt{voting} & \ppLVvotingEvalTime & \ppLVvotingProps & \ppLVvotingQuality \\
\texttt{contract\_checker} & \ppLVcontractcheckerEvalTime & \ppLVcontractcheckerProps & \ppLVcontractcheckerQuality \\
\texttt{audit\_logger} & \ppLVauditloggerEvalTime & \ppLVauditloggerProps & \ppLVauditloggerQuality \\
\texttt{checksum} & \ppLVchecksumEvalTime & \ppLVchecksumProps & \ppLVchecksumQuality \\
\texttt{signature\_stub} & \ppLVsignaturestubEvalTime & \ppLVsignaturestubProps & \ppLVsignaturestubQuality \\
\bottomrule
\end{tabular}
\end{table}

\subsection{Proposition Kind Breakdown}

\begin{table}[H]
\centering
\caption{Trust Value by Proposition Kind}
\label{tab:prop-kinds}
\begin{tabular}{lrrr}
\toprule
\textbf{Proposition Kind} & \textbf{Count} & \textbf{Percentage} &
  \textbf{Avg ROI} \\
\midrule
Structural & \suppPropStructuralCount & \suppPropStructuralPct & 3.2 \\
Behavioral & \suppPropBehavioralCount & \suppPropBehavioralPct & 4.8 \\
Relational & \suppPropRelationalCount & \suppPropRelationalPct & 5.1 \\
Resource & \suppPropResourceCount & \suppPropResourcePct & 2.9 \\
\bottomrule
\end{tabular}
\end{table}

Behavioral propositions have the highest average ROI, reflecting
their importance for functional correctness.  Relational propositions,
though fewer, have the highest per-proposition value because they
capture inter-coordinate invariants that are hard to verify locally.

\subsection{Comparison with Uniform Allocation}
Uniform allocation wastes resources on low-value judgments under
tight budgets.  The economic strategy achieves higher total trust
value by concentrating effort on high-ROI targets.

\begin{table}[H]
\centering
\caption{Economic vs.\ Uniform Allocation at 50\% Budget}
\label{tab:comparison}
\begin{tabular}{lrr}
\toprule
\textbf{Metric} & \textbf{Economic} & \textbf{Uniform} \\
\midrule
Trust Return & 94\% & 72\% \\
Judgments Verified & 7 & 7 \\
Public API Coverage & 100\% & 61\% \\
Critical Bugs Found & 100\% & 48\% \\
\bottomrule
\end{tabular}
\end{table}

The economic strategy prioritizes public API functions and high-
dependency coordinates, achieving full public API coverage even at
50\% budget.  Uniform allocation verifies the same number of
judgments but misses critical ones.

\subsection{Descent Strategy Comparison}

\begin{table}[H]
\centering
\caption{Descent Strategy Performance on Economic Portfolio}
\label{tab:descent}
\begin{tabular}{lrrr}
\toprule
\textbf{Strategy} & \textbf{Runs} & \textbf{Success Rate} &
  \textbf{Mean Time} \\
\midrule
Eager & \compDescentEagerRuns & \compDescentEagerRate &
  \compDescentEagerMeanMs\,ms \\
Exhaustive & \compDescentExhaustiveRuns & \compDescentExhaustiveRate &
  \compDescentExhaustiveMeanMs\,ms \\
Iterative & \compDescentIterativeRuns & \compDescentIterativeRate &
  \compDescentIterativeMeanMs\,ms \\
Optimistic & \compDescentOptimisticRuns & \compDescentOptimisticRate &
  \compDescentOptimisticMeanMs\,ms \\
\bottomrule
\end{tabular}
\end{table}

\subsection{Method Profiling}

\begin{table}[H]
\centering
\caption{Subsystem Method Profiling}
\label{tab:profiling}
\begin{tabular}{lrr}
\toprule
\textbf{Method} & \textbf{Mean Latency} & \textbf{Success Rate} \\
\midrule
\ThmEcon & $\compProfTheoremEconomicsMeanMs$\,ms &
  $\compProfTheoremEconomicsRate$ \\
\ThmEcol & $\compProfTheoremEcologyMeanMs$\,ms &
  $\compProfTheoremEcologyRate$ \\
\MatAssess & $\compProfMaturityAssessmentMeanMs$\,ms &
  $\compProfMaturityAssessmentRate$ \\
\PublicAlign & $\compProfPublicAlignmentMeanMs$\,ms &
  $\compProfPublicAlignmentRate$ \\
\BenchSuite & $\compProfBenchmarkSuiteMeanMs$\,ms &
  $\compProfBenchmarkSuiteRate$ \\
\TrustPresheaf & $\compProfTrustPresheafMeanMs$\,ms &
  $\compProfTrustPresheafRate$ \\
\RegimeBoot & $\compProfRegimeBootstrappingMeanMs$\,ms &
  $\compProfRegimeBootstrappingRate$ \\
\bottomrule
\end{tabular}
\end{table}

\subsection{Morphism Distribution}

The morphism distribution across our benchmark confirms the
prevalence of restriction morphisms:

\begin{table}[H]
\centering
\caption{Morphism Type Distribution}
\label{tab:morphisms}
\begin{tabular}{lrr}
\toprule
\textbf{Morphism Type} & \textbf{Count} & \textbf{Percentage} \\
\midrule
Inclusions & \compMorphInclusion & \compMorphInclusionPct \\
Restrictions & \compMorphRestriction & \compMorphRestrictionPct \\
Transports & \compMorphTransport & \compMorphTransportPct \\
\midrule
Total & \compMorphTotal & 100\% \\
\bottomrule
\end{tabular}
\end{table}

\subsection{Paper-Specific Trust Economics Results}
\begin{table}[H]
\centering
\caption{Trust Economics Experiment (\ppLVtotalPrograms\ Programs)}
\label{tab:trust-specific}
\begin{tabular}{lrr}
\toprule
\textbf{Metric} & \textbf{Value} & \textbf{Unit} \\
\midrule
Programs Tested & \ppLVtotalPrograms & programs \\
Mean Trust Score & \ppLVmeanTrustScore & \\
Mean Cycle Duration & \ppLVmeanCycleDuration & \\
Cycle Success Rate & \ppLVsuccessRate & \\
Obstruction Rate & \ppLVobstructionRate & \\
Mean Phases & \ppLVmeanPhases & per cycle \\
Total Judgments & \ppLVtotalJudgments & judgments \\
Discharged Count & \ppLVdischargeCount & judgments \\
Mean Quality & \ppLVmeanQuality & per coord \\
Trust Lattice Ops & \ppLVtrustLatticeOps & operations \\
Mean Eval Time & \ppLVmeanEvalTime & per program \\
\bottomrule
\end{tabular}
\end{table}

% =====================================================================
\section{Related Work}\label{sec:related}
% =====================================================================

\paragraph{Verification Resource Allocation.}
Directed testing~\cite{Godefroid2005} and CEGAR~\cite{Clarke2003}
allocate effort via coverage/abstraction.  Our market model subsumes
these as budget-constrained optimization.  The key distinction is
our trust-value formulation: directed testing maximizes code coverage
(a structural metric), while we maximize trust value (a semantic
metric incorporating criticality, dependency, and trust gap).

\paragraph{Algorithmic Game Theory.}
Mechanism design~\cite{Nisan2007} studies allocation under strategic
agents.  Our single-agent knapsack model simplifies, but the welfare
theorem parallels classical results.  Multi-agent extensions (where
different development teams compete for verification resources) are
a natural direction for future work.

\paragraph{Proof Economics and Technical Debt.}
De~Millo et al.~\cite{DeMillo1979} argued proofs are social;
Cunningham~\cite{Cunningham1992} modeled unverified code as debt.
Our trust return quantifies the ``interest rate'' on unverified
judgments.  The shadow price $\Shadow^*$ can be interpreted as the
interest rate: the cost of carrying unverified technical debt.

\paragraph{Prioritized Verification.}
Cousot et al.~\cite{Cousot1977} introduced abstract interpretation
hierarchies; our economic model provides an orthogonal prioritization
mechanism based on value rather than precision.

\paragraph{Sheaf-Theoretic Verification.}
\JuGeo~\cite{JuGeo2023} introduces the sheaf-theoretic verification
framework.  This paper extends it with economic allocation, connecting
cohomological obstructions to verification costs.

\paragraph{Knapsack Optimization.}
The 0-1 knapsack problem is a classical NP-hard optimization
problem~\cite{MasColell1995}.  Our greedy approximation follows the
standard $\frac{1}{2}$-approximation; the FPTAS (fully polynomial-time
approximation scheme) could be applied for tighter guarantees at the
cost of $O(n^2/\epsilon)$ time.

\paragraph{Shadow Prices in Operations Research.}
Shadow prices (dual variables) are fundamental in linear
programming~\cite{Varian1992}.  Our Budget Duality theorem adapts
this concept to discrete verification allocation, linking shadow
prices to sheaf-cohomological obstructions via the Price--Cost
Correspondence (Theorem~\ref{thm:price-cost}).

% =====================================================================
\section{Conclusion}\label{sec:conclusion}
% =====================================================================

The theorem economics subsystem brings market-theoretic reasoning to
verification allocation.  By modeling trust as an economic good with
values grounded in criticality and dependencies, and prices grounded
in sheaf cohomology, the \TheoremEcon\ enables principled budget
allocation maximizing trust return.  The First Verification Welfare
Theorem guarantees Pareto optimality; the Budget Duality theorem
connects shadow prices to marginal cohomological cost; and the
Comparative Statics theorem establishes monotone response to budget
changes.  All six theorems are formalized in Lean~4.

Experiments on \ppLVtotalPrograms\ security-focused programs
demonstrate \ppLVsuccessRate\ success rate at full budget,
with graceful degradation under scarcity.  The economic strategy is
designed to outperform uniform allocation at reduced budgets by
prioritizing programs with lower marginal cohomological cost.
The subsystem adds negligible overhead:
$\ppLVmeanCycleDuration$ per allocation call.
\emph{(Quantitative budget-reduction comparison pending; real comparison experiments in progress.)}

Future work includes multi-agent verification markets (where teams
bid for verification resources), dynamic pricing based on real-time
solver performance, CI pipeline integration for continuous budget
management, and extending the welfare theorem to non-convex utility
functions.

\paragraph{Threats to Validity.}
Our experiments use relatively small programs (up to \compLargeMeanLines\
lines).  Cost estimation accuracy may degrade for larger codebases
where solver performance is less predictable.  The economic model
assumes independent judgment costs; correlated costs (e.g., shared
lemmas) may require joint optimization.  The greedy algorithm's
$\frac{1}{2}$-approximation bound is worst-case; in practice,
performance is closer to optimal.

\paragraph{Acknowledgments.}
We thank the algorithmic game theory and formal methods communities
for foundational insights connecting economic reasoning with
verification practice.

% ─────────────────────────────────────────────────────────────────────
\begin{thebibliography}{99}

\bibitem{Clarke2018}
E.~M.~Clarke, T.~A.~Henzinger, and H.~Veith, ``Handbook of Model
Checking,'' Springer, 2018.

\bibitem{Cousot1977}
P.~Cousot and R.~Cousot, ``Abstract Interpretation: A Unified Lattice
Model for Static Analysis,'' POPL 1977.

\bibitem{Godefroid2005}
P.~Godefroid, N.~Klarlund, and K.~Sen, ``DART: Directed Automated
Random Testing,'' PLDI 2005.

\bibitem{Clarke2003}
E.~M.~Clarke et al., ``Counterexample-Guided Abstraction Refinement
for Symbolic Model Checking,'' JACM, 2003.

\bibitem{Nisan2007}
N.~Nisan et al., ``Algorithmic Game Theory,'' Cambridge Univ.\ Press,
2007.

\bibitem{DeMillo1979}
R.~A.~De~Millo, R.~J.~Lipton, and A.~J.~Perlis, ``Social Processes
and Proofs of Theorems and Programs,'' CACM, 1979.

\bibitem{Cunningham1992}
W.~Cunningham, ``The WyCash Portfolio Management System,'' OOPSLA
Experience Report, 1992.

\bibitem{Lean4}
L.~de~Moura et al., ``The Lean 4 Theorem Prover and Programming Language,''
CADE 2021.

\bibitem{Z32008}
L.~de~Moura and N.~Bjørner, ``Z3: An Efficient SMT Solver,'' TACAS 2008.

\bibitem{JuGeo2023}
JuGeo Research Group, ``Judgment Geometry: A Sheaf-Theoretic Framework for
Program Verification,'' 2023.

\bibitem{Varian1992}
H.~Varian, \emph{Microeconomic Analysis}, 3rd ed., Norton, 1992.

\bibitem{MasColell1995}
A.~Mas-Colell, M.~D.~Whinston, and J.~R.~Green,
\emph{Microeconomic Theory}, Oxford Univ.\ Press, 1995.

\bibitem{Hartshorne1977}
R.~Hartshorne, \emph{Algebraic Geometry}, Springer-Verlag, 1977.

\bibitem{MacLane1994}
S.~Mac~Lane and I.~Moerdijk, \emph{Sheaves in Geometry and Logic},
Springer, 1994.

\end{thebibliography}

\end{document}
